% !TeX program = xelatex

\documentclass[11pt]{article}


%\usepackage[a4paper,margin=1in]{geometry}
\usepackage[a4paper, left=0.5in, right=2in, top=1in, bottom=1in, marginparwidth=1.2in]{geometry}
\usepackage[T1]{fontenc}
\usepackage[utf8]{inputenc}
% lmodern removed for compile compat

\usepackage{amsmath,amsthm,amssymb,amsfonts}
\usepackage{hyperref}
\usepackage{enumitem}
\usepackage[protrusion=true,expansion=false]{microtype}

\hypersetup{
  colorlinks=true,
  linkcolor=blue,
  urlcolor=blue,
  citecolor=blue
}

\newcounter{ReviewerNoteCount}
\setcounter{ReviewerNoteCount}{0}

\usepackage{xcolor} % 引入颜色包用于高亮

\newcommand{\reviewer}[2]{%
    \stepcounter{ReviewerNoteCount}% 每次调用时计数器自动加 1
    {\color{blue}#1}% 原文标蓝（保留原文内容）
    \textsuperscript{\color{blue}[\theReviewerNoteCount]}% 在原文后添加红色上标，如 [1]
    \marginpar{\linespread{1.1}\selectfont\color{red}\scriptsize\raggedright\textbf{[\theReviewerNoteCount]} #2}% 侧边栏显示带编号的批注
}

% 如果您发现页边距不够宽导致批注显示不全，可以在 \geometry 中适当增加 right margin，
% 例如将 \usepackage[a4paper,margin=1in]{geometry} 改为：
% \usepackage[a4paper, left=1in, right=1.5in, top=1in, bottom=1in, marginparwidth=1.2in]{geometry}
% --- 审阅人批注自定义命令结束 ---


\numberwithin{equation}{section}

\newtheorem{theorem}{Theorem}[section]
\newtheorem{lemma}[theorem]{Lemma}
\newtheorem{proposition}[theorem]{Proposition}
\newtheorem{corollary}[theorem]{Corollary}
\theoremstyle{definition}
\newtheorem{definition}[theorem]{Definition}
\newtheorem{remark}[theorem]{Remark}
\newtheorem{example}[theorem]{Example}

% -------------------------------------------------------

\title{\textbf{Projection, Measure, and Idempotent Relations:\\
Independent Axioms and a Fixed-Point Coupling Law}}
\author{
Yunbeom Yi\\[4pt]
\small Independent Researcher, Republic of Korea\\
\small Email: \texttt{28ho2132@gmail.com}
}
\date{}

% -------------------------------------------------------

\begin{document}
\maketitle

\begin{abstract}
We introduce a minimal \reviewer{ZFC--internal}{ZFC-internal} axiom system for \emph{pre-structural data}
\[
(X,\mathcal A,\mu,\mu^{\otimes2},R,I,\Pi_R,G,E_0,\eta),
\]
where $\Pi_R:X\to R$ is a designated map and $G\subseteq X\times X$ is a measurable relation. Those pre-structural data satisfying Axioms~I--III are called \emph{admissible structural models}; the axioms couple a finitely additive measure, an \reviewer{idempotent retraction, and an idempotent symmetric}{idempotent retraction, and an idempotent, symmetric} relation through a single coupling law (Axiom~III). The main results are the following. \reviewer{Consistency: the}{Consistency: The} axiom system is satisfiable in ZFC, with explicit finite and countable models, including finite families for which $\eta\neq 0$. \reviewer{Independence: the}{Independence: The} three axioms, and the three subclauses of Axiom~III (conservation, \reviewer{projection--measure}{projection-measure} invariance, coupling law), are mutually independent. The coupling law admits a \reviewer{fixed--point}{fixed-point} reformulation: \reviewer{it is the unique}{It is the unique} bounded finitely additive solution of a \reviewer{Banach--contraction}{Banach contraction (standard fixed term, no hyphen)} equation $f=T_\eta f$ determined by $(\mu,\Pi_R,\eta)$, with closed form $f_*(B)=\mu(B)+\tfrac{\eta}{1-\eta}\mu(\Pi_R^{-1}(B))$. Admissible structural models with fixed $\eta$ form a category $\mathbf{Struct}_\eta$; under fiber measurability and either a \reviewer{finiteness~(R--fin)}{finiteness (R-fin)} or countability plus \reviewer{$\sigma$--additivity}{$\sigma$-additivity}~(R--ctbl) hypothesis, a \reviewer{quotient--factorization}{quotient-factorization} theorem reduces the classification problem to the \reviewer{identity--retraction}{identity-retraction} case $\Pi_R=\mathrm{id}_X$, where each \reviewer{$G$--equivalence}{$G$-equivalence} class $C_k$ satisfies $\mu(C_k)\in\{0,(1-\eta)^{-1}\}$ under the hypotheses of Theorem~\ref{thm:pi-id-classification}.
\end{abstract}




\bigskip

\noindent\textbf{Keywords:}
Axiomatic Systems; Set Theory (ZFC); Measure Theory; Charges and Finitely Additive Measures; Idempotent Relations; Retraction Operators; Coupling Law; Banach Fixed-Point Theorem; Independence of Axioms; Categorical Structure; Quotient Factorization.

% -------------------------------------------------------

\section{Introduction}

We present a small \reviewer{ZFC--internal}{ZFC-internal} axiom system for abstract \emph{pre-structural data}, tuples
\[
\mathcal D=(X,\mathcal A,\mu,\mu^{\otimes2},R,I,\Pi_R,G,E_0,\eta)
\]
consisting of:
(i) a \reviewer{set algebra}{algebra of sets (consistent with Definition 2.1)} $(X,\mathcal A)$ with a finitely additive measure $\mu:\mathcal A\to[0,\infty)$ and a finitely additive product charge $\mu^{\otimes2}$ on $\mathcal A\otimes\mathcal A$ satisfying the rectangle rule;
(ii) disjoint measurable sets $R,I\in\mathcal A$ and a map $\Pi_R:X\to R$ (not yet required to be idempotent);
(iii) a measurable relation $G\subseteq X\times X$ (not yet required to be reflexive, symmetric, or idempotent);
(iv) two scalars $E_0>0$ and $\eta\in[0,1]$.

Three axioms tie these components together. Axiom~I imposes projective structure (idempotence of $\Pi_R$ and $\Pi_R|_R=\mathrm{id}_R$, plus measurability of preimages of measurable subsets of $R$). Axiom~II imposes closure structure on $G$ (reflexivity, symmetry, and the idempotence $G\circ G=G$). Axiom~III imposes quantitative coupling between measure, projection, and relation via three subclauses: conservation ($\mu(R)+\mu(I)=E_0>0$), projection--measure invariance ($\mu(\Pi_R^{-1}(B))=\mu(B)$ for measurable $B\subseteq R$), and a coupling law on $(B\times X)\cap G$. A pre-structural datum satisfying Axioms~I--III is called an \emph{admissible structural model}.

Our goal is to isolate a minimal \reviewer{first--order}{first-order} framework in which projection, closure, and measure coexist and interact.
\emph{We do not assume $X=R\sqcup I$ at the axiomatic level;} only $R\cap I=\varnothing$ and measurability of $R\cup I$ are required. When the full partition $X=R\sqcup I$ is additionally imposed, the invariance clause of Axiom~III forces $\mu(I)=0$ (Proposition~\ref{prop:muI-zero}).
Despite their minimality, the axioms yield explicit finite and countable admissible models as well as several structural consequences.
We use finite additivity to keep the axiom system minimal; Section~\ref{sec:sigma-additive} records the additional constraints under \reviewer{$\sigma$--additivity}{$\sigma$-additivity}.

\paragraph{On terminology.}
We adopt the name \emph{curvature} for the idempotent symmetric relation $G$ and for the derived \reviewer{two--point}{two-point} quantity $\mu^{\otimes2}((B\times X)\cap G)$ as a convenient label for a bilinear, \reviewer{self--composing}{self-composing} interaction pattern. This usage is analogical rather than \reviewer{differential--geometric}{differential-geometric}: no Riemannian, sectional, or Ricci curvature is defined or used in this paper. The terminology emphasizes that $G\circ G=G$ imposes a \reviewer{closure--type}{closure-type} rigidity on pair interactions, in the same way that curvature operators in geometry control how infinitesimal data compose. The reader for whom this analogy is unhelpful may replace ``curvature'' with ``idempotent interaction'' throughout without loss of content.

\subsection{Motivation}

Projection, closure, and measure are standard objects studied in different areas of mathematics (see e.g.\ \cite{Halmos,Bogachev,AliprantisBorder} for measure theory, \cite{Rao,Fremlin} for charges, \cite{Kuratowski} for topology, \cite{BauschkeCombettes} for\reviewer{projection/retraction}{projection and retraction} operators in analysis, \cite{KolokoltsovMaslov} for idempotent analysis, \cite{KrantzEtAl} for axiomatic foundations of measurement, and \cite{MacLane,Awodey,Riehl} for categorical viewpoints).
The axioms place them into a single framework with direct interaction:
\[
\mu^{\otimes2}((B\times X)\cap G)
= \mu(B) + \eta\,\mu^{\otimes2}((\Pi_R^{-1}(B)\times X)\cap G).
\]
This coupling law does not appear in standard \reviewer{measure--theoretic}{measure-theoretic} or reviewer{closure--theoretic}{closure-theoretic} axiomatizations.
It couples the measure, projection, and closure components into a single axiom.

\subsection{Relation to classical frameworks}

Each component has a classical analogue:
\begin{itemize}[nosep]
  \item $G\circ G = G$ encodes an idempotent closure/\reviewer{equivalence--type}{equivalence-type} behavior at the level of relations.
  \item $\Pi_R\circ\Pi_R = \Pi_R$ is the defining idempotence of a retraction.
  \item $(X,\mathcal A,\mu)$ is a finitely additive measure space.
  \item The invariance $\mu(\Pi_R^{-1}(B))=\mu(B)$ for $B\subseteq R$ is analogous to a \reviewer{measure--preserving}{measure-preserving} retraction onto $R$.
\end{itemize}
The key point is that these components coexist in one structure and are coupled by Axiom~III.

\subsection{Structural unification}

The axioms provide a common interface for several classical viewpoints:
\begin{itemize}[nosep]
  \item \textbf{Measure/Probability}: if one normalizes $\mu(X)=1$, the measure component can be read as a finitely additive probability.
  \item \textbf{Closure/Topology}: $G$ induces a closure/connectedness pattern via its idempotence.
  \item \textbf{Equivalence}: on each \reviewer{$G$--connected}{$G$-connected} component, idempotence forces transitivity along \reviewer{$G$--chains}{$G$-chains}.
  \item \textbf{Projection}: $\Pi_R$ preserves $\mu$ on measurable subsets of $R$ (Axiom~III).
\end{itemize}
Thus, measure, projection, and closure appear as compatible components of a single \reviewer{ZFC--internal}{ZFC-internal} structure.

\subsection{Contributions}\label{sec:contributions}
The contributions can be summarized as follows.
\begin{itemize}[nosep]
  \item A \reviewer{ZFC--internal}{ZFC-internal} axiom system (Section~\ref{sec:axioms}) and explicit model families establishing consistency (Section~\ref{sec:models}), including finite families with $\eta\neq 0$.
  \item Independence of the three axioms and of the three subclauses of Axiom~III, by separating models (Theorem~\ref{thm:independence-main}, Corollary~\ref{cor:minimality}).
  \item A \reviewer{fixed--point}{fixed-point} reformulation of the coupling law as the \reviewer{Banach--contraction}{Banach contraction} equation $f=T_\eta f$, with closed form $f_*(B)=\mu(B)+\tfrac{\eta}{1-\eta}\mu(\Pi_R^{-1}(B))$ and a \reviewer{Neumann--series}{Neumann series} expansion (Theorem~\ref{thm:coupling-fixed-point}, Corollary~\ref{cor:coupling-identification}).
  \item A \reviewer{quotient--factorization}{quotient-factorization} theorem reducing admissibility to the \reviewer{identity--retraction}{identity-retraction} case $\Pi_R=\mathrm{id}_X$ under fiber measurability and either (R--fin) or (R--ctbl) (Theorem~\ref{thm:quotient-factorization}). In that case, under the hypotheses of Theorem~\ref{thm:pi-id-classification}, each \reviewer{$G$--equivalence}{$G$-equivalence} class carries mass $0$ or $(1-\eta)^{-1}$.
\end{itemize}
Section~\ref{sec:sigma-additive} records a side observation: under \reviewer{$\sigma$--additivity}{$\sigma$-additivity} with $\mu(X)=1$, the global instance of the coupling law forces $\eta=0$; so nontrivial $\eta\neq 0$ families appear either in the finitely additive setting or for \reviewer{$\sigma$--additive}{$\sigma$-additive} measures not normalized as probabilities.


\section{Language and basic objects}

We recall standard notions used throughout the paper.

\begin{definition}[Algebra of sets]
Let $X$ be nonempty. A family $\mathcal A\subseteq\mathcal P(X)$ is an \emph{algebra} if $\varnothing,X\in\mathcal A$ and
\reviewer{$B_1,B_2\in\mathcal A$ implies}{$B_1,B_2\in\mathcal A$ imply}
$B_1\cup B_2\in\mathcal A$ and $X\setminus B_1\in\mathcal A$.
\end{definition}

\begin{definition}[Finitely additive measure]
A finitely additive measure is a map $\mu:\mathcal A\to[0,\infty)$ such that $\mu(\varnothing)=0$ and
$\mu(B_1\sqcup B_2)=\mu(B_1)+\mu(B_2)$ whenever $B_1,B_2\in\mathcal A$ are disjoint.
\end{definition}

\begin{definition}[Binary relations and composition]
A binary relation $H$ on $X$ is a subset of $X\times X$. \reviewer{Its composition with a relation $K$ is}{The composition of $H$ with a relation $K$ is}
\[
H\circ K = \{(x,z):\exists y,\ (x,y)\in H,\ (y,z)\in K\}.
\]
The diagonal is $\Delta_X=\{(x,x):x\in X\}$. A relation is \emph{idempotent} if $H\circ H=H$.
\end{definition}

\begin{definition}[Retraction]
If $R\subseteq X$, a map $\Pi_R:X\to R$ is a \emph{retraction} if $\Pi_R|_R=\mathrm{id}_R$.
\reviewer{We also impose idempotence}{A retraction is additionally required to be idempotent} $\Pi_R\circ\Pi_R=\Pi_R$.
\end{definition}

\begin{definition}[Product algebra and product charge]\label{def:product-measure}
Given $(X,\mathcal A,\mu)$, the product algebra $\mathcal A\otimes\mathcal A$ is the algebra generated by rectangles
$B_1\times B_2$ with $B_i\in\mathcal A$.

A finitely additive set function $\mu^{\otimes2}:\mathcal A\otimes\mathcal A\to[0,\infty)$ is called a \emph{product charge}
(for $\mu$) if it satisfies the rectangle rule
\[
\mu^{\otimes2}(B_1\times B_2)=\mu(B_1)\mu(B_2)\qquad(B_1,B_2\in\mathcal A).
\]
In the finitely additive setting, \reviewer{existence/uniqueness}{existence and uniqueness} of such an extension is not automatic in full generality; whenever
$\mu^{\otimes2}$ is used below, it is treated as part of the model data satisfying the rectangle rule.
(Under \reviewer{$\sigma$--additivity}{$\sigma$-additivity} and \reviewer{$\sigma$--finiteness}{$\sigma$-finiteness}, the standard product charge exists uniquely; see Section~\ref{sec:sigma-additive}.)
\end{definition}

\begin{lemma}[Atomic product charge exists]\label{lem:atomic-product-charge}
Assume $X$ is \reviewer{finite or countable}{finite and countable} and $\mathcal A=\mathcal P(X)$. Suppose $\mu$ is atomic, i.e.\ there exists
$m:X\to[0,\infty)$ such that $\mu(B)=\sum_{x\in B}m(x)$ for all $B\subseteq X$ (with the sum finite in the finite case).
Define $\mu^{\otimes2}:\mathcal P(X\times X)\to[0,\infty)$ by
\[
\mu^{\otimes2}(S):=\sum_{(x,y)\in S} m(x)m(y)\qquad(S\subseteq X\times X).
\]
Then $\mu^{\otimes2}$ is finitely additive and satisfies the rectangle rule
$\mu^{\otimes2}(A\times B)=\mu(A)\mu(B)$ for all $A,B\subseteq X$.
In particular, in the \reviewer{finite/countable}{finite and countable} atomic regimes used in Section~5, a product charge can be taken canonically.
\end{lemma}

\begin{proof}
Finite additivity follows from disjointness of sums. For rectangles,
\[
\mu^{\otimes2}(A\times B)=\sum_{x\in A}\sum_{y\in B} m(x)m(y)
=\Big(\sum_{x\in A}m(x)\Big)\Big(\sum_{y\in B}m(y)\Big)=\mu(A)\mu(B).
\]
\end{proof}


\section{The axioms}\label{sec:axioms}

\begin{definition}[Pre-structural datum]\label{def:structural-datum}
A \emph{pre-structural datum} is a tuple
\[
\mathcal D=(X,\allowbreak \mathcal A,\allowbreak \mu,\allowbreak \mu^{\otimes2},\allowbreak R,\allowbreak I,\allowbreak \Pi_R,\allowbreak G,\allowbreak E_0,\allowbreak \eta)
\]
such that:
\begin{itemize}[nosep]
  \item $X$ is a nonempty set and $\mathcal A\subseteq\mathcal P(X)$ is an algebra;
  \item $\mu:\mathcal A\to[0,\infty)$ is a finitely additive measure;
  \item $\mu^{\otimes2}:\mathcal A\otimes\mathcal A\to[0,\infty)$ is a finitely additive product charge satisfying the rectangle rule
  \[
  \mu^{\otimes2}(B_1\times B_2)=\mu(B_1)\mu(B_2)\qquad(B_1,B_2\in\mathcal A);
  \]
  \item $R,I\in\mathcal A$ are disjoint and $R\cup I\in\mathcal A$;
  \item $\Pi_R:X\to R$ is a map;
  \item $G\subseteq X\times X$ lies in $\mathcal A\otimes\mathcal A$;
  \item $E_0\in(0,\infty)$ and $\eta\in[0,1]$ are scalars.
\end{itemize}
The pre-structural datum specifies only the raw signature of a structural model; the structural constraints (idempotence of $\Pi_R$, properties of $G$, compatibility with $\mu$) are imposed separately by Axioms~I--III below.
\end{definition}

\paragraph{Axiom I (duality and projection).}
$R\cap I=\varnothing$ and $R\cup I\in\mathcal A$.
$\Pi_R$ is idempotent and $\Pi_R|_R=\mathrm{id}_R$.
If $B\in\mathcal A\cap\mathcal P(R)$, then $\Pi_R^{-1}(B)\in\mathcal A$.

\paragraph{Axiom II (closed curvature).}
$G$ is reflexive, symmetric, and idempotent: $G\circ G=G$.

\paragraph{Axiom III (Coupling).}
\reviewer{$\mu(R)+\mu(I)=E_0>0$.
If $B\subseteq R$ is measurable then $\mu(\Pi_R^{-1}(B))=\mu(B)$.
For all $B\in\mathcal A$,
\[
\mu^{\otimes2}((B\times X)\cap G)
= \mu(B)+\eta\,\mu^{\otimes2}((\Pi_R^{-1}(B)\times X)\cap G).
\]}{Theorem 3's clause (c) is cited more than once in the following text, but it is not numbered here.}

\begin{remark}\label{rem:preimage}
Since $\Pi_R$ maps into $R$, one has $\Pi_R^{-1}(B)=\Pi_R^{-1}(B\cap R)$ for all $B\in\mathcal A$.
Thus $\Pi_R^{-1}(B)$ is measurable whenever $B\cap R$ is, by Axiom~I.
\end{remark}

\begin{definition}[Admissible structural model]\label{def:admissible}
A pre-structural datum
\[
\mathcal D=(X,\mathcal A,\mu,\mu^{\otimes2},R,I,\Pi_R,G,E_0,\eta)
\]
is called an \emph{admissible structural model} (or simply a \emph{structural model}) if it satisfies Axioms~I, II, and~III. When we refer to $\mathcal M$ as a \emph{structural model} without further qualification throughout the paper, admissibility is understood unless explicitly stated otherwise. The objects of the categories $\mathbf{Struct}_\eta$ and $\mathbf{Struct}^{\mathrm{id}}_\eta$ introduced in Section~\ref{sec:category} are admissible structural models by convention.
\end{definition}


\section{Basic consequences}

We record consequences of the axioms.

\begin{proposition}[Elementary consequences of Axioms I and II]\label{prop:elementary}
From Axioms~I and~II one has immediately:
\begin{enumerate}[label=(\roman*),nosep]
\item (\emph{retraction fixed points}) $\Pi_R(x)=x$ if and only if $x\in R$;
\item (\emph{iterated projection}) $\Pi_R^n(x)=\Pi_R(x)$ for every integer $n\ge 1$; in particular $\Pi_R^n(x)\in R$, so \reviewer{a point not initially in $I$ is never sent into $I$}{a point not in $I$ is never mapped into $I$} by any iterate of $\Pi_R$;
\item (\emph{$G$ is an equivalence relation}) \reviewer{from reflexivity, symmetry}{From reflexivity, symmetry}, and idempotence $G\circ G=G$, transitivity follows: \reviewer{if $(x,y),(y,z)\in G$ then}{if $(x,y),(y,z)\in G$, then} $(x,z)\in G\circ G=G$.
\end{enumerate}
\end{proposition}

\begin{proof}
\reviewer{(i) follows from}{(i) Follows from} $\Pi_R(X)\subseteq R$ and $\Pi_R|_R=\mathrm{id}_R$ (Axiom~I). \reviewer{(ii) is a consequence}{(ii) Is a consequence} of the idempotence $\Pi_R\circ\Pi_R=\Pi_R$: for $n=1$ this is definitional, and for $n\ge 2$ we have $\Pi_R^n(x)=\Pi_R(\Pi_R^{n-1}(x))=\Pi_R(x)$ by induction. \reviewer{(iii) is immediate}{(iii) Is immediate} from Axiom~II as stated.
\end{proof}

\begin{proposition}[Projection--measure invariance]\label{prop:projection-measure}
If $B\in\mathcal A\cap\mathcal P(R)$, then $\mu(\Pi_R^{-1}(B))=\mu(B)$.
\end{proposition}

\begin{proof}
This is the second assertion of Axiom~III.
\end{proof}

\begin{proposition}[Support concentration under full partition]\label{prop:muI-zero}
Assume in addition that $X=R\sqcup I$. Then Axiom~III (projection--measure invariance) forces $\mu(I)=0$.
\end{proposition}

\begin{proof}
Take $B=R$ in $\mu(\Pi_R^{-1}(B))=\mu(B)$ for measurable $B\subseteq R$.
Since $\Pi_R^{-1}(R)=X$, we have $\mu(X)=\mu(R)$.
By finite additivity on $X=R\sqcup I$, $\mu(X)=\mu(R)+\mu(I)$, hence $\mu(I)=0$.
\end{proof}

\begin{remark}[Observational weight and hidden sector]
Proposition~\ref{prop:muI-zero} should be read as a design feature rather than a degeneracy: $\mu$ encodes
\emph{observational weight} carried by the representative sector $R$.
The hidden sector $I$ may be structurally present (e.g.\ as fibers of $\Pi_R$ and through the interaction channel $G$),
while being forced to carry no \reviewer{$\mu$--mass}{$\mu$-mass} under full-partition invariance.
In particular, the axioms do \emph{not} assume $X=R\sqcup I$; the full-partition hypothesis is an optional strengthening.
\end{remark}

\begin{example}[Why invariance forces zero mass on nontrivial fibers]\label{ex:fiber-mass}
Let $X=\{r_0,r_1,a,b\}$, $R=\{r_0,r_1\}$, and $I=\{a,b\}$, and define a retraction $\Pi_R:X\to R$ by
$\Pi_R|_R=\mathrm{id}_R$ together with $\Pi_R(a)=r_0$ and $\Pi_R(b)=r_1$.
If $\mu(\{a\})>0$, then for $B=\{r_0\}\subseteq R$ one has $\Pi_R^{-1}(B)=\{r_0,a\}$ and hence
\[
\mu(\Pi_R^{-1}(B))=\mu(\{r_0\})+\mu(\{a\})>\mu(\{r_0\})=\mu(B),
\]
so the invariance $\mu(\Pi_R^{-1}(B))=\mu(B)$ fails.
Thus, under Axiom~III, positive mass cannot be carried on nontrivial fibers of $\Pi_R$.
\end{example}

\begin{corollary}[Baseline case $\eta=0$]\label{cor:eta0-baseline}
If $\eta=0$, the coupling law of Axiom~III reduces to
\[
\mu^{\otimes2}((B\times X)\cap G)=\mu(B)\qquad(\forall\,B\in\mathcal A),
\]
so the observed curvature load coincides with the base measure.
\end{corollary}

\begin{proof}
Set $\eta=0$ in the coupling law of Axiom~III.
\end{proof}


\subsection{Fixed-point characterization of the coupling law}\label{sec:fixed-point-char}

We record an \reviewer{operator--theoretic}{operator-theoretic} reformulation of \reviewer{subclause~(c) of Axiom~III}{subclause (c) of Axiom III} (the coupling law) as a \reviewer{fixed--point}{fixed-point} equation determining a unique bounded extension of $\mu$ under the projection $\Pi_R$.

Throughout this subsection we assume Axioms~I and~II hold, and that $\mu$ is bounded: $\mu(X)\le M<\infty$. Let
\[
\mathcal B:=\{\,f:\mathcal A\to\mathbb{R}\ \mid\ f\ \text{finitely additive and bounded}\,\}
\]
denote the space of bounded finitely additive signed charges, equipped with the supremum norm
\[
\|f\|_\infty:=\sup_{B\in\mathcal A}|f(B)|.
\]
Under this norm $\mathcal B$ is a Banach space (see, e.g., \cite{Rao,AliprantisBorder}). The subset $\mathcal B_+$ of nonnegative $f$ is a closed convex cone in $\mathcal B$.

\begin{definition}[Projection--update operator]\label{def:T-eta}
For $\eta\in[0,1)$, define the operator $T_\eta:\mathcal B\to\mathcal B$ by
\[
(T_\eta f)(B):=\mu(B)+\eta\,f(\Pi_R^{-1}(B))\qquad(B\in\mathcal A).
\]
\end{definition}

That $T_\eta$ is \reviewer{well--defined}{well-defined} into $\mathcal B$ follows from three observations: (i) $\Pi_R^{-1}(B)\in\mathcal A$ for all $B\in\mathcal A$ (Remark~\ref{rem:preimage}); (ii) both $\mu$ and $f\circ\Pi_R^{-1}$ are finitely additive in $B$, hence so is $T_\eta f$; (iii) if $\|f\|_\infty<\infty$ then $|T_\eta f(B)|\le M+\eta\|f\|_\infty$, so $T_\eta f\in\mathcal B$. Moreover $T_\eta$ preserves nonnegativity: if $f\ge 0$ on $\mathcal A$ and $\mu\ge 0$, then $T_\eta f\ge 0$, so $T_\eta(\mathcal B_+)\subseteq\mathcal B_+$.

\begin{theorem}[Coupling law as unique fixed point]\label{thm:coupling-fixed-point}
Assume Axioms~I and~II, and $\mu(X)\le M<\infty$. For $\eta\in[0,1)$, the equation
\[
f=T_\eta f
\]
has a unique solution $f_*\in\mathcal B$, given explicitly by
\[
f_*(B)=\mu(B)+\frac{\eta}{1-\eta}\,\mu(\Pi_R^{-1}(B))\qquad(B\in\mathcal A).
\]
Moreover, $f_*\in\mathcal B_+$, and for any starting point $f_0\in\mathcal B$ the iterates $f_n:=T_\eta^n f_0$ satisfy
\[
\|f_n-f_*\|_\infty\le\eta^n\|f_0-f_*\|_\infty,
\]
so $f_n\to f_*$ in norm.
\end{theorem}

\begin{proof}
\emph{Contraction.} For $f_1,f_2\in\mathcal B$ and any $B\in\mathcal A$,
\[
|(T_\eta f_1)(B)-(T_\eta f_2)(B)|=\eta\,|f_1(\Pi_R^{-1}(B))-f_2(\Pi_R^{-1}(B))|\le\eta\|f_1-f_2\|_\infty.
\]
Taking the supremum over $B$ yields $\|T_\eta f_1-T_\eta f_2\|_\infty\le\eta\|f_1-f_2\|_\infty$. Since $\eta<1$, $T_\eta$ is a strict contraction on the Banach space $\mathcal B$. By Banach's \reviewer{fixed--point}{fixed-point} theorem there exists a unique $f_*\in\mathcal B$ with $f_*=T_\eta f_*$, together with the stated geometric convergence rate for iterates.

\emph{Explicit form.} We verify that the stated closed form $f_*(B)=\mu(B)+\tfrac{\eta}{1-\eta}\mu(\Pi_R^{-1}(B))$ satisfies $f_*=T_\eta f_*$. The key identity is the \emph{idempotence of the preimage map}:
\[
\Pi_R^{-1}(\Pi_R^{-1}(B))=\Pi_R^{-1}(B)\qquad(\forall B\in\mathcal A).
\]
Indeed, for any $x\in X$, $x\in\Pi_R^{-1}(\Pi_R^{-1}(B))$ if and only if $\Pi_R(x)\in\Pi_R^{-1}(B)$, if and only if $\Pi_R(\Pi_R(x))\in B$, if and only if $\Pi_R(x)\in B$ (by idempotence of $\Pi_R$ from Axiom~I, see Proposition~\ref{prop:elementary}(ii)), if and only if $x\in\Pi_R^{-1}(B)$. Hence
\begin{align*}
(T_\eta f_*)(B)
&=\mu(B)+\eta\,f_*(\Pi_R^{-1}(B))\\
&=\mu(B)+\eta\Bigl[\mu(\Pi_R^{-1}(B))+\tfrac{\eta}{1-\eta}\,\mu\bigl(\Pi_R^{-1}(\Pi_R^{-1}(B))\bigr)\Bigr]\\
&=\mu(B)+\eta\,\mu(\Pi_R^{-1}(B))+\tfrac{\eta^2}{1-\eta}\,\mu(\Pi_R^{-1}(B))\\
&=\mu(B)+\Bigl[\eta+\tfrac{\eta^2}{1-\eta}\Bigr]\mu(\Pi_R^{-1}(B))\\
&=\mu(B)+\tfrac{\eta}{1-\eta}\,\mu(\Pi_R^{-1}(B))=f_*(B).
\end{align*}
Hence $f_*$ is a fixed point; by uniqueness, it is the fixed point. Nonnegativity $f_*\ge 0$ is immediate since $\mu\ge 0$ and $\eta/(1-\eta)\ge 0$.
\end{proof}

\begin{corollary}[Coupling law as identification of curvature load]\label{cor:coupling-identification}
Assume Axioms~I and~II, and $\mu(X)<\infty$. For $\eta\in[0,1)$, the coupling law (subclause~(c) of Axiom~III) is equivalent to the identification of the observed curvature load with the unique fixed point of $T_\eta$:
\[
\mu^{\otimes2}\bigl((B\times X)\cap G\bigr)=f_*(B)=\mu(B)+\frac{\eta}{1-\eta}\,\mu(\Pi_R^{-1}(B))
\qquad(\forall B\in\mathcal A).
\]
\end{corollary}

\begin{proof}
Let $\mathcal C(B):=\mu^{\otimes2}((B\times X)\cap G)$. The coupling law asserts $\mathcal C(B)=\mu(B)+\eta\,\mathcal C(\Pi_R^{-1}(B))$ for all $B\in\mathcal A$, i.e.\ $\mathcal C=T_\eta\mathcal C$. Finite additivity of $\mathcal C$ in $B$ follows from the rectangle rule and finite additivity of $\mu^{\otimes2}$, applied to the disjoint decomposition $(B_1\sqcup B_2)\times X=(B_1\times X)\sqcup(B_2\times X)$. Boundedness: $\mathcal C(B)\le\mu^{\otimes2}(X\times X)=\mu(X)^2<\infty$. Hence $\mathcal C\in\mathcal B$. By Theorem~\ref{thm:coupling-fixed-point} the only fixed point of $T_\eta$ in $\mathcal B$ is $f_*$, so $\mathcal C=f_*$.
\end{proof}

\begin{remark}[Two compatibility conditions determining the coupling law]\label{rem:coupling-not-ad-hoc}
Corollary~\ref{cor:coupling-identification} expresses the coupling law as the unique bounded finitely additive solution of the operator equation $f=T_\eta f$. The operator $T_\eta$ encodes two conditions on a quantity $f$ defined on $\mathcal A$:
\begin{itemize}[nosep]
  \item \emph{Baseline agreement with $\mu$}: for $\eta=0$ one has $T_0 f=\mu$ and $f_*=\mu$.
  \item \emph{\reviewer{Self--similar}{Self-similar} propagation along $\Pi_R$}: the correction to $\mu$ propagates under the preimage map $\Pi_R^{-1}$, scaled by $\eta$.
\end{itemize}
Requiring the curvature load $\mathcal C(B)=\mu^{\otimes2}((B\times X)\cap G)$ to satisfy both conditions determines the coupling law uniquely once $(\mathcal A,\mu,\Pi_R,\eta)$ are fixed.
\end{remark}

\begin{remark}[Neumann--series expansion]\label{rem:neumann}
Formally iterating $T_\eta$ from any $f_0\in\mathcal B$ and using Banach convergence, one obtains the \reviewer{Neumann--series}{Neumann series} expression
\[
f_*(B)=\sum_{n=0}^{\infty}\eta^n\,\mu(\Pi_R^{-n}(B)),
\]
convergent in $\|\cdot\|_\infty$ for $\eta\in[0,1)$. Using the preimage idempotence $\Pi_R^{-n}(B)=\Pi_R^{-1}(B)$ for $n\ge 1$ (derived in the proof of Theorem~\ref{thm:coupling-fixed-point}), the series telescopes to
\[
f_*(B)=\mu(B)+\mu(\Pi_R^{-1}(B))\sum_{n=1}^{\infty}\eta^n=\mu(B)+\frac{\eta}{1-\eta}\,\mu(\Pi_R^{-1}(B)),
\]
recovering the closed form. The term $\mu(B)$ is the baseline ($n=0$) contribution; the geometric sum $\tfrac{\eta}{1-\eta}$ accumulates all iterated projection corrections.
\end{remark}

\begin{remark}[Relation to the classification when $\Pi_R=\mathrm{id}$]\label{rem:relation-classification}
When $\Pi_R=\mathrm{id}_X$, one has $\Pi_R^{-1}(B)=B$, so the \reviewer{fixed--point}{fixed-point} formula collapses to
\[
f_*(B)=\frac{\mu(B)}{1-\eta}\qquad(B\ \text{measurable}).
\]
In regimes where the full hypotheses of Theorem~\ref{thm:pi-id-classification} hold for the relevant \reviewer{identity--retraction}{identity-retraction} core, this yields the blockwise normalization law: every \reviewer{positive--mass}{positive-mass} \reviewer{$G$--equivalence}{$G$-equivalence} class $C_k$ satisfies $\mu(C_k)=(1-\eta)^{-1}$.
\end{remark}


\subsection{Observable reduction under global decoupling}

\begin{definition}[Observed curvature load]\label{def:observed-curvature-load}
For measurable $B\subseteq R$, define
\[
\mathcal C_R(B):=\mu^{\otimes2}\big((B\times X)\cap G\big).
\]
\end{definition}

\begin{theorem}[Exact reduction on $R$ under global decoupling]\label{thm:observable-reduction}
Let $\mathcal M=(X,\allowbreak \mathcal A,\allowbreak \mu,\allowbreak \mu^{\otimes2},\allowbreak R,\allowbreak I,\allowbreak \Pi_R,\allowbreak G,\allowbreak E_0,\allowbreak \eta)$ be a structural model with $\eta\in[0,1)$.
Assume the global decoupling condition
\[
\mu^{\otimes2}\big(((\Pi_R^{-1}(B)\setminus B)\times X)\cap G\big)=0
\qquad(\forall\,B\subseteq R\ \text{measurable}).
\]
Then for every measurable $B\subseteq R$:
\begin{enumerate}[label=(\alph*),nosep]
\item (\emph{support reduction}) $\mu^{\otimes2}\big(G\setminus(R\times R)\big)=0$ and
\[
\mu^{\otimes2}\big((B\times X)\cap G\big)=\mu^{\otimes2}\big((B\times R)\cap(G\cap(R\times R))\big).
\]
\item (\emph{closed form on $R$}) the coupling law closes on $R$:
\[
\mu^{\otimes2}\big((B\times X)\cap G\big)=\frac{\mu(B)}{1-\eta}.
\]
\end{enumerate}
\end{theorem}

\begin{proof}
Fix measurable $B\subseteq R$.

\reviewer{(\emph{support reduction.})}{(a) Support reduction.}
Apply the decoupling condition with $B=R$. Since $\Pi_R^{-1}(R)=X$, we have $\Pi_R^{-1}(R)\setminus R=X\setminus R$, hence
\[
0=\mu^{\otimes2}\big(((X\setminus R)\times X)\cap G\big).
\]
By symmetry of $G$, also $\mu^{\otimes2}((X\times(X\setminus R))\cap G)=0$, and therefore
$\mu^{\otimes2}(G\setminus(R\times R))=0$.
In particular, $(B\times X)\cap G$ and $(B\times R)\cap(G\cap(R\times R))$ differ by a \reviewer{$\mu^{\otimes2}$--null}{$\mu^{\otimes2}$-null} set, yielding the stated equality.

\reviewer{(\emph{closed form.})}{(b) Closed form.}
Since $B\subseteq \Pi_R^{-1}(B)$ and $\Pi_R^{-1}(B)=B\sqcup(\Pi_R^{-1}(B)\setminus B)$, we have a disjoint union
\[
(\Pi_R^{-1}(B)\times X)\cap G
=
((B\times X)\cap G)\ \sqcup\ \big(((\Pi_R^{-1}(B)\setminus B)\times X)\cap G\big).
\]
By the decoupling hypothesis, the second term has \reviewer{$\mu^{\otimes2}$--measure}{$\mu^{\otimes2}$-measure} zero, hence
\[
\mu^{\otimes2}\big((\Pi_R^{-1}(B)\times X)\cap G\big)=\mu^{\otimes2}\big((B\times X)\cap G\big).
\]
Insert this into the coupling law (Axiom~III) for $B\subseteq R$:
\[
\mu^{\otimes2}\big((B\times X)\cap G\big)=\mu(B)+\eta\,\mu^{\otimes2}\big((\Pi_R^{-1}(B)\times X)\cap G\big),
\]
to obtain $(1-\eta)\mu^{\otimes2}((B\times X)\cap G)=\mu(B)$, i.e.\ the stated closed form.
\end{proof}


\section{Models and consistency}\label{sec:models}

We construct explicit models.

\subsection{Finite models}

\begin{theorem}\label{thm:finite-model}
There exists a finite structural model. In particular, the axiom set is satisfiable in ZFC.
\end{theorem}


  
  

\begin{proof}
Let $X$ be finite and nonempty, and choose disjoint nonempty subsets $R,I\subseteq X$
Let $X$ be finite and nonempty, and choose disjoint nonempty subsets $R,I\subseteq X$
with $R\cup I = X$.
Let $\mathcal A=\mathcal P(X)$.
Fix $r_0\in R$ and define $\Pi_R:X\to R$ by
\[
\Pi_R(x)=
\begin{cases}
x, & x\in R,\\
r_0, & x\in I.
\end{cases}
\]
Then $\Pi_R$ is idempotent and $\Pi_R|_R=\mathrm{id}_R$.
Since $\mathcal A$ is the full power set, $\Pi_R^{-1}(B)\in\mathcal A$ for all $B\subseteq R$,
so Axiom~I holds.

Define $G=\Delta_X\subseteq X\times X$.
Then $G$ is reflexive, symmetric, and idempotent; also $G\in\mathcal A\otimes\mathcal A$.
Thus Axiom~II holds.

Choose weights $m(x)\in\{0,1\}$ such that $m(x)=0$ for all $x\in I$ and
$\sum_{x\in R} m(x)>0$.
Define $\mu:\mathcal A\to[0,\infty)$ by
\[
\mu(B)=\sum_{x\in B} m(x).
\]
Then $\mu$ is finitely additive, $\mu(I)=0$, and $\mu(R)>0$.
Set
\[
E_0 := \mu(X) = \mu(R) \in (0,\infty).
\]
This gives the conservation identity $\mu(R)+\mu(I)=E_0>0$.

If $B\subseteq R$, then
\[
\Pi_R^{-1}(B) =
\begin{cases}
B, & r_0\notin B,\\
B\cup I, & r_0\in B,
\end{cases}
\]
so
\[
\mu(\Pi_R^{-1}(B))
=
\begin{cases}
\mu(B), & r_0\notin B,\\
\mu(B)+\mu(I), & r_0\in B,
\end{cases}
=
\mu(B),
\]
because $\mu(I)=0$.
Thus the projection--measure invariance part of Axiom~III holds.

Using the product charge $\mu^{\otimes2}$,
For any $B\subseteq X$,
\[
(B\times X)\cap G
= \{(x,x): x\in B\},
\]
and therefore
\[
\mu^{\otimes2}((B\times X)\cap G)
= \sum_{x\in B} \mu(\{x\})^2
= \sum_{x\in B} m(x)^2
= \sum_{x\in B} m(x)
= \mu(B),
\]
using $m(x)^2=m(x)$ for $m(x)\in\{0,1\}$.
Taking $\eta=0$, the coupling law in Axiom~III reduces to
$\mu^{\otimes2}((B\times X)\cap G)=\mu(B)$, which we have just verified.
Thus we obtain a finite structural model satisfying all axioms.
\end{proof}


\subsection{Further finite models with \texorpdfstring{$\eta\neq 0$}{eta ≠ 0}}

The preceding constructions set $\eta=0$ for simplicity. We now record explicit finite models
and families in which $\eta\neq 0$.

\begin{theorem}[Finite model with $\eta\neq 0$]\label{thm:eta-nonzero-finite}
Let $X=\{r_0,r_1,i\}$ and define a measurable partition $X=R\sqcup I$ by
\[
R=\{r_0,r_1\},\qquad I=\{i\}.
\]
Define a projection (retraction) $\Pi_R:X\to R$ by
\[
\Pi_R(x)=
\begin{cases}
x,& x\in R,\\
r_0,& x=i\in I.
\end{cases}
\]
Let $G=\Delta_X=\{(x,x):x\in X\}\subseteq X\times X$. Fix any $\eta\in[0,1)$ with $\eta\neq 0$.
Define a finitely additive measure $\mu$ on $\mathcal P(X)$ by point masses
\[
m(i)=0,\qquad m(r_0)=m(r_1)=\frac{1}{1-\eta},\qquad \mu(B):=\sum_{x\in B} m(x).
\]
Define the product charge $\mu^{\otimes 2}$ on $\mathcal P(X\times X)$ by rectangle masses
$\mu^{\otimes 2}(A\times B):=\mu(A)\mu(B)$ and finite additivity.

Then $\mathcal M=(X,\allowbreak \mathcal P(X),\allowbreak \mu,\allowbreak \mu^{\otimes2},\allowbreak R,\allowbreak I,\allowbreak \Pi_R,\allowbreak G,\allowbreak E_0,\allowbreak \eta)$ satisfies Axioms I--III, with
$E_0=\mu(R)+\mu(I)=\frac{2}{1-\eta}>0$.
\end{theorem}

\begin{proof}
Axiom~I and Axiom~II are immediate from the definitions. For Axiom~III, conservation holds by the choice of $E_0$.
If $B\subseteq R$, then $\Pi_R^{-1}(B)$ equals $B$ or $B\cup\{i\}$ depending on whether $r_0\in B$, and $m(i)=0$ yields
$\mu(\Pi_R^{-1}(B))=\mu(B)$.

For the coupling law, note that $(B\times X)\cap \Delta_X=\{(x,x):x\in B\}$, so
\[
\mu^{\otimes 2}((B\times X)\cap G)=\sum_{x\in B} m(x)^2.
\]
Similarly, $\mu^{\otimes 2}((\Pi_R^{-1}(B)\times X)\cap G)=\sum_{x\in \Pi_R^{-1}(B)} m(x)^2$, and the extra term (if any) is $m(i)^2=0$.
Thus the coupling law reduces to checking for each $r\in R$:
\[
m(r)^2=m(r)+\eta\,m(r)^2 \iff (1-\eta)m(r)^2=m(r),
\]
which holds for $m(r)=\frac{1}{1-\eta}$.
\end{proof}

\begin{proposition}[A finite model with $\eta\neq 0$ and non-diagonal curvature]\label{prop:eta-nonzero-nondiag}
Let $X=\{r_0,r_1\}$ and take $R=X$, $I=\varnothing$, $\Pi_R=\mathrm{id}_X$.
Fix any $\eta\in[0,1)$, $\eta\neq 0$. Define $\mu$ by
\[
m(r_0)=m(r_1)=\frac{1}{2(1-\eta)},\qquad \mu(B)=\sum_{x\in B}m(x),
\]
so that $\mu(X)=\frac{1}{1-\eta}$. Let $G=X\times X$ (the total relation on $X$).
Let $\mu^{\otimes2}$ be the canonical atomic product charge on $\mathcal P(X\times X)$ induced by $\mu$ (Lemma~\ref{lem:atomic-product-charge}).
Then $\mathcal M=(X,\allowbreak \mathcal P(X),\allowbreak \mu,\allowbreak \mu^{\otimes2},\allowbreak R,\allowbreak I,\allowbreak \Pi_R,\allowbreak G,\allowbreak E_0,\allowbreak \eta)$ satisfies Axioms I--III with $E_0=\mu(R)=\frac{1}{1-\eta}$, and $G$ is non-diagonal.
\end{proposition}

\begin{proof}
Invariance is immediate since $\Pi_R=\mathrm{id}_X$. For any $B\subseteq X$, $(B\times X)\cap G=B\times X$, hence
\[
\mu^{\otimes 2}((B\times X)\cap G)=\mu(B)\mu(X).
\]
Also $\Pi_R^{-1}(B)=B$, so the coupling law becomes $\mu(B)\mu(X)=\mu(B)+\eta\,\mu(B)\mu(X)$.
If $\mu(B)>0$ we cancel $\mu(B)$ to obtain $(1-\eta)\mu(X)=1$, which holds by construction; if $\mu(B)=0$ the identity is trivial.
\end{proof}
\begin{theorem}[Canonical classification when $\Pi_R=\mathrm{id}$]\label{thm:pi-id-classification}
Assume $R=X$, $I=\varnothing$, and $\Pi_R=\mathrm{id}_X$. Let $\eta\in[0,1)$.
Let $G\subseteq X\times X$ be reflexive, symmetric, and idempotent ($G\circ G=G$).
Then $G$ is an equivalence relation on $X$; denote its equivalence classes by $(C_k)_{k\in K}$, so that
\[
G=\bigcup_{k\in K}(C_k\times C_k)\qquad\text{(disjoint union of blocks)}.
\]
Let $\mu$ be a finitely additive measure on $\mathcal P(X)$ and let $\mu^{\otimes2}$ be the product charge on the product algebra $\mathcal P(X)\otimes\mathcal P(X)$, satisfying the rectangle rule of Definition~\ref{def:product-measure}. Suppose further that one of the following summability hypotheses holds, so that the block formula below is well--defined:
\begin{itemize}[nosep]
\item[\emph{(K--fin)}] the index set $K$ is finite; or
\item[\emph{(K--ctbl)}] $K$ is countable, $\mu$ is $\sigma$--additive (so $\mathcal P(X)$ is treated as a $\sigma$--algebra) and $\sigma$--finite on $X$, and $\mu^{\otimes2}$ is the standard $\sigma$--additive product measure on the product $\sigma$--algebra $\mathcal P(X)\otimes\mathcal P(X)$ (which then exists and is uniquely characterized by the rectangle rule).
\end{itemize}
Then the coupling law
\[
\mu^{\otimes2}((B\times X)\cap G)=\mu(B)+\eta\,\mu^{\otimes2}((B\times X)\cap G)\qquad(\forall\,B\subseteq X\ \text{measurable})
\]
holds if and only if, for every class $C_k$ with $\mu(C_k)>0$,
\[
\mu(C_k)=\frac{1}{1-\eta}.
\]
Equivalently, for all measurable $B\subseteq X$ one has the explicit block formula
\[
\mu^{\otimes2}((B\times X)\cap G)=\sum_{k\in K}\mu(B\cap C_k)\,\mu(C_k),
\]
and the coupling law is precisely the blockwise normalization law above.
\end{theorem}

\begin{proof}
First, $G$ is transitive: if $(x,y)\in G$ and $(y,z)\in G$, then $(x,z)\in G\circ G=G$.
With reflexivity and symmetry, $G$ is an equivalence relation, hence decomposes into blocks $(C_k)$ and
$G=\bigcup_k (C_k\times C_k)$.

For any measurable $B\subseteq X$,
\[
(B\times X)\cap G=\bigsqcup_{k\in K}\big((B\cap C_k)\times C_k\big),
\]
a disjoint union. Under~(K--fin) finite additivity of $\mu^{\otimes2}$ together with the rectangle rule gives
\[
\mu^{\otimes2}((B\times X)\cap G)=\sum_{k\in K}\mu(B\cap C_k)\,\mu(C_k)
\]
immediately. Under~(K--ctbl) the same identity follows from $\sigma$--additivity of $\mu^{\otimes2}$ applied to the countable disjoint union together with the rectangle rule. In either case we obtain the stated block formula.

Since $\Pi_R=\mathrm{id}_X$, the coupling law reduces to
\[
(1-\eta)\,\mu^{\otimes2}((B\times X)\cap G)=\mu(B)\qquad(\forall\,B\subseteq X\ \text{measurable}).
\]
Taking $B=C_j$ yields $(1-\eta)\mu(C_j)^2=\mu(C_j)$, hence $\mu(C_j)=0$ or $\mu(C_j)=\frac{1}{1-\eta}$.
Thus all positive-mass classes satisfy the normalization $\mu(C_k)=\frac{1}{1-\eta}$.
Conversely, if this holds for all positive-mass classes, then the block formula implies
$\mu^{\otimes2}((B\times X)\cap G)=\mu(B)/(1-\eta)$ for all measurable $B$, so the coupling law holds.
\end{proof}

\begin{remark}[On the summability hypothesis]\label{rem:block-summability}
The block formula of Theorem~\ref{thm:pi-id-classification} involves a sum over $K$, which can fail to make sense in the purely finitely additive setting when $K$ is infinite: arbitrary countable disjoint unions need not be countably additive for a finitely additive $\mu^{\otimes2}$. The hypotheses~(K--fin) or~(K--ctbl) are exactly what guarantees the block formula as an identity; under neither hypothesis, the theorem holds only in the weaker form ``for each finite sub-collection of classes, the corresponding partial block formula holds'', which is insufficient for the global coupling--law characterization. We therefore always state downstream results (including Corollary~\ref{cor:curvature-load-closed-form} and Theorem~\ref{thm:quotient-factorization}) under these summability hypotheses when the block formula is invoked.
\end{remark}

\begin{corollary}[Total mass in the fully positive block case]\label{cor:pi-id-total-mass}
Under the hypotheses of Theorem~\ref{thm:pi-id-classification}, suppose that $X$ has finitely many
$G$--equivalence classes $C_1,\dots,C_m$ and that $\mu(C_j)>0$ for all $j$.
Then
\[
\mu(X)=\frac{m}{1-\eta}
\qquad\text{and}\qquad
\mu^{\otimes2}(G)=\sum_{j=1}^m \mu(C_j)^2=\frac{m}{(1-\eta)^2}.
\]
Moreover, for every $B\subseteq X$ one has the closed form
\[
\mu^{\otimes2}((B\times X)\cap G)=\frac{\mu(B)}{1-\eta}.
\]
\end{corollary}

\begin{proof}
By Theorem~\ref{thm:pi-id-classification}, $\mu(C_j)=\frac{1}{1-\eta}$ for each $j$.
Since $X=\bigsqcup_{j=1}^m C_j$, finite additivity gives $\mu(X)=\sum_{j=1}^m \mu(C_j)=\frac{m}{1-\eta}$.
Also $G=\bigcup_{j=1}^m (C_j\times C_j)$ is a disjoint union, hence
\[
\mu^{\otimes2}(G)=\sum_{j=1}^m \mu(C_j)\mu(C_j)=\sum_{j=1}^m \mu(C_j)^2=\frac{m}{(1-\eta)^2}.
\]
Finally, the same theorem yields $\mu^{\otimes2}((B\times X)\cap G)=\mu(B)/(1-\eta)$ for all $B\subseteq X$.
\end{proof}

\begin{remark}
Theorem~\ref{thm:pi-id-classification} subsumes the block--curvature family (Theorem~\ref{thm:block-curvature-family}) by deriving the block decomposition directly from the axioms on $G$.
\end{remark}

\begin{theorem}[Block-curvature family for $\eta\neq 0$]\label{thm:block-curvature-family}
Let $X$ be finite and let $\{C_1,\dots,C_m\}$ be a partition of $X$ into nonempty blocks.
Set $R=X$, $I=\varnothing$, $\Pi_R=\mathrm{id}_X$, and define
\[
G:=\bigcup_{k=1}^m (C_k\times C_k)\subseteq X\times X.
\]
Fix $\eta\in[0,1)$ with $\eta\neq 0$. Let $\mu$ be any finitely additive measure on $\mathcal P(X)$ and define $\mu^{\otimes 2}$ on rectangles by
$\mu^{\otimes 2}(A\times B)=\mu(A)\mu(B)$.

Then the coupling law holds for all $B\subseteq X$ if and only if
\[
\mu(C_k)=\frac{1}{1-\eta}\qquad \text{for all }k=1,\dots,m.
\]
In particular, at fixed $\eta\neq 0$ the coupling law imposes exactly $m$ scalar constraints (one per $G$-component), and imposes no further constraint on the distribution of $\mu$ inside each block.
\end{theorem}

\begin{proof}
For any $B\subseteq X$,
\[
(B\times X)\cap G=\bigcup_{k=1}^m \big((B\cap C_k)\times C_k\big)
\]
is a disjoint union, hence
\[
\mu^{\otimes 2}((B\times X)\cap G)=\sum_{k=1}^m \mu(B\cap C_k)\,\mu(C_k).
\]
Since $\Pi_R^{-1}(B)=B$, the coupling law is equivalent to
\[
(1-\eta)\sum_{k=1}^m \mu(B\cap C_k)\,\mu(C_k)=\mu(B)\qquad(\forall B\subseteq X).
\]
Taking $B=C_j$ yields $(1-\eta)\mu(C_j)^2=\mu(C_j)$, so $\mu(C_j)=0$ or $\mu(C_j)=\frac{1}{1-\eta}$.
For nontrivial blocks we take the positive solution and obtain necessity.
Conversely, if $\mu(C_k)=\frac{1}{1-\eta}$ for all $k$, then
\[
\mu^{\otimes 2}((B\times X)\cap G)=\frac{1}{1-\eta}\sum_{k=1}^m \mu(B\cap C_k)=\frac{1}{1-\eta}\mu(B),
\]
which satisfies the coupling law identically.
\end{proof}

\begin{remark}[Interpretation: blockwise normalization law]\label{rem:block-norm}
For $\Pi_R=\mathrm{id}$, the coupling law acts as a \emph{blockwise normalization law}:
it fixes the total mass of each $G$--connected component to $(1-\eta)^{-1}$, while leaving internal distributions free.
The endpoint $\eta=1$ would force $\mu(C_k)=+\infty$ and is therefore excluded in these finite families.
\end{remark}


We next give a finite model with nontrivial curvature.

\begin{proposition}[Finite class model with nontrivial curvature]\label{prop:finite-class-model}
Let $X$ be a finite nonempty set, partitioned into nonempty equivalence classes
$(C_j)_{j\in J}$.
For each $j$, choose a distinguished representative $r_j\in C_j$ and set
\[
R := \{r_j : j\in J\},\qquad I := X\setminus R.
\]
Let $\mathcal A=\mathcal P(X)$.
Choose weights $w_j\in\{0,1\}$ with $\sum_j w_j>0$ and define a finitely additive measure
$\mu:\mathcal A\to[0,\infty)$ by
\[
\mu(B) := \sum_{j\in J : r_j\in B} w_j.
\]
Define $\Pi_R:X\to R$ by $\Pi_R(x)=r_j$ whenever $x\in C_j$, and define
\[
G := \bigcup_{j\in J} (C_j\times C_j)\subseteq X\times X.
\]
Set $E_0 := \mu(X)$ and $\eta:=0$.
Let $\mu^{\otimes2}$ be the canonical atomic product charge on $\mathcal P(X\times X)$ induced by $\mu$ (Lemma~\ref{lem:atomic-product-charge}).
Then
\[
\mathcal M=(X,\mathcal A,\mu,\mu^{\otimes2},R,I,\Pi_R,G,E_0,\eta)
\]
is a structural model.  In particular, there exist finite models with nontrivial curvature.
\end{proposition}

\begin{proof}
By construction, $R,I\in\mathcal A$, $R\cap I=\varnothing$, and $R\cup I=X$.
The map $\Pi_R$ restricts to the identity on $R$, and $\Pi_R\circ\Pi_R=\Pi_R$ because
every $x$ is sent to its class representative.
Preimages of subsets of $R$ are unions of classes, hence measurable.
Thus Axiom~I holds.

The relation $G$ is reflexive, symmetric, and idempotent since it is the union of
equivalence relations on the blocks $C_j$.
Clearly $G\in\mathcal A\otimes\mathcal A$, so Axiom~II holds.

For Axiom~III, note first that
\[
\mu(R)+\mu(I) = \mu(R) = \sum_j w_j =:E_0>0.
\]
If $B\subseteq R$, then
\[
\Pi_R^{-1}(B) = \bigcup_{j : r_j\in B} C_j,
\]
and therefore
\[
\mu(\Pi_R^{-1}(B))
= \sum_{j : r_j\in B} \mu(C_j)
= \sum_{j : r_j\in B} w_j
= \mu(B).
\]
Thus the projection--measure invariance holds.

Equip $X\times X$ with the product charge $\mu^{\otimes2}$.
For any $B\subseteq X$, write $B_j:=B\cap C_j$.
Then
\[
(B\times X)\cap G
= \bigcup_{j\in J} \bigl( B_j\times C_j \bigr),
\]
disjoint union.
Hence
\[
\mu^{\otimes2}((B\times X)\cap G)
= \sum_{j\in J} \mu(B_j)\,\mu(C_j)
= \sum_{j : r_j\in B} w_j^2
= \sum_{j : r_j\in B} w_j
= \mu(B),
\]
using $\mu(B_j)=w_j$ if $r_j\in B_j$ and $0$ otherwise, and $w_j^2=w_j$ for $w_j\in\{0,1\}$.
With $\eta=0$, the coupling law reduces to $\mu^{\otimes2}((B\times X)\cap G)=\mu(B)$.
Thus $\mathcal M$ satisfies all axioms and has nontrivial curvature whenever some class
$C_j$ has cardinality $>1$.
\end{proof}

\subsection{Countable model}

\begin{proposition}\label{prop:countable-model}
There exists a structural model with $X$ countably infinite.
\end{proposition}

\begin{proof}
Let $\mathcal M_0=(X_0,\mathcal A_0,\mu_0,\mu_0^{\otimes2},R_0,I_0,\Pi_{R,0},G_0,E_{0,0},\eta_0)$
be a finite structural model as in Theorem~\ref{thm:finite-model}, with $\eta_0=0$.
Choose a countably infinite set $Y$ disjoint from $X_0$ and set
\[
X := X_0 \sqcup Y.
\]
Define $\mathcal A$ to be the algebra generated by $\mathcal A_0$ together with all subsets of $Y$.
More concretely, $\mathcal A$ consists of all sets of the form
\[
B_0 \sqcup B_Y,\qquad B_0\in\mathcal A_0,\ B_Y\subseteq Y.
\]

Define $\mu:\mathcal A\to[0,\infty)$ by
\[
\mu(B_0\sqcup B_Y) := \mu_0(B_0),
\]
so that all points of $Y$ have measure zero.
Set
\[
R := R_0,\qquad
I := I_0 \sqcup Y.
\]
Define $\Pi_R:X\to R$ by
\[
\Pi_R(x) :=
\begin{cases}
\Pi_{R,0}(x), & x\in X_0,\\
r_0, & x\in Y,
\end{cases}
\]
where $r_0\in R_0$ is the distinguished point used in the finite model.
Define $G\subseteq X\times X$ by
\[
G := G_0 \cup \Delta_Y,
\]
where $\Delta_Y=\{(y,y):y\in Y\}$.
Set $E_0 := \mu(R)+\mu(I) = \mu_0(R_0)+\mu_0(I_0) = E_{0,0}$ and $\eta := 0$.

We verify the axioms.
Axiom~I holds since $R,I\in\mathcal A$, $R\cap I=\varnothing$, $R\cup I=X$,
$\Pi_R$ is idempotent and equals $\Pi_{R,0}$ on $R_0$.
Measurability of preimages for subsets of $R$ follows from the definition of $\mathcal A$.

Axiom~II holds because $G_0$ and $\Delta_Y$ are both reflexive, symmetric, and idempotent,
and their union is block--diagonal with respect to the decomposition $X=X_0\sqcup Y$.
Moreover, $G\in\mathcal A\otimes\mathcal A$.

For Axiom~III, note first that
\[
\mu(R)+\mu(I) = \mu_0(R_0)+\mu_0(I_0) = E_{0,0} > 0.
\]
If $B\subseteq R$ is measurable, then $B\subseteq R_0$ and
\[
\Pi_R^{-1}(B) = \Pi_{R,0}^{-1}(B) \sqcup (Y\ \text{if } r_0\in B\ \text{else }\varnothing),
\]
so
\[
\mu(\Pi_R^{-1}(B))
= \mu_0(\Pi_{R,0}^{-1}(B)) = \mu_0(B) = \mu(B),
\]
using the finite model's Axiom~III.
Thus the projection--measure invariance holds.

Equip $X\times X$ with the product charge $\mu^{\otimes2}$.
If $B\in\mathcal A$ and we write $B=B_0\sqcup B_Y$ with
$B_0\in\mathcal A_0$, $B_Y\subseteq Y$, then
\[
(B\times X)\cap G
= \bigl( (B_0\times X_0)\cap G_0 \bigr)
  \ \sqcup\ \{(y,y):y\in B_Y\},
\]
and
\[
\mu^{\otimes2}((B\times X)\cap G)
= \mu_0^{\otimes2}((B_0\times X_0)\cap G_0)
  + \sum_{y\in B_Y} \mu(\{y\})^2
= \mu_0^{\otimes2}((B_0\times X_0)\cap G_0),
\]
since $\mu(\{y\})=0$ for $y\in Y$.
Similarly,
\[
\mu(B) = \mu_0(B_0).
\]
With $\eta=0$, Axiom~III reduces to
\[
\mu^{\otimes2}((B\times X)\cap G) = \mu(B),
\]
which we have just verified using the corresponding identity in the finite model.
Thus $\mathcal M=(X,\mathcal A,\mu,\mu^{\otimes2},R,I,\Pi_R,G,E_0,\eta)$ is a countably infinite structural model.
\end{proof}

% -------------------------------------------------------

\section{Independence of the axioms}\label{sec:independence}

In this section we establish that the three axioms of Section~\ref{sec:axioms} are mutually independent: no axiom follows from the remaining two. Together with the model constructions of Section~\ref{sec:models}, which establish consistency, this shows that the axiom system is both consistent and minimal, justifying the word ``minimal'' used in the abstract and Section~\ref{sec:contributions}.

Independence is established in the model-theoretic sense: for each axiom $\mathsf X\in\{\mathrm{I},\mathrm{II},\mathrm{III}\}$ we exhibit a structure over the same signature which satisfies the other two axioms but fails $\mathsf X$. We also record the stronger result that the three subclauses of Axiom~III (conservation, projection--measure invariance, coupling law) are themselves mutually independent within Axiom~III.

\subsection{Separating models for Axioms I, II, III}

The constructions below follow a minimal relaxation principle: each separating datum $\mathcal D_{\neg \mathsf X}$ is obtained from a baseline finite model (as in Section~\ref{sec:models}) by modifying exactly the feature asserted by the axiom under analysis, leaving the remaining structural features intact.

\begin{proposition}[Independence of Axiom I]\label{prop:indep-I}
There exists a pre-structural datum $\mathcal D_{\neg \mathrm{I}}$ over the ambient signature of Definition~\ref{def:structural-datum} which satisfies Axioms~II and~III but fails Axiom~I.
\end{proposition}

\begin{proof}
Let $X=\{a,b,c\}$ and $\mathcal A=\mathcal P(X)$. Set
\[
R:=\{a,b\},\qquad I:=\{c\},
\]
so $R\cap I=\varnothing$ and $R\cup I=X\in\mathcal A$. Define the map $\Pi_R:X\to R$ by
\[
\Pi_R(a):=b,\qquad \Pi_R(b):=a,\qquad \Pi_R(c):=a.
\]
Then $\Pi_R(X)\subseteq R$, but $\Pi_R|_R\ne\mathrm{id}_R$ (it swaps $a$ and $b$) and $\Pi_R\circ\Pi_R\ne\Pi_R$ (since $\Pi_R(\Pi_R(a))=\Pi_R(b)=a\ne b=\Pi_R(a)$). Hence Axiom~I fails.

Set $G:=\Delta_X$, which is reflexive, symmetric, and idempotent on any set; hence Axiom~II holds.

Define $\mu:\mathcal A\to[0,\infty)$ by the point weights $m(a)=m(b):=1$ and $m(c):=0$, and let $\mu^{\otimes2}$ be the canonical atomic product charge from Lemma~\ref{lem:atomic-product-charge}. Set $\eta:=0$ and $E_0:=\mu(R)+\mu(I)=2>0$, so Axiom~III(a) (conservation) holds.

For projection--measure invariance, take any $B\subseteq R$ measurable:
\[
\Pi_R^{-1}(\{a\})=\{b,c\},\quad \Pi_R^{-1}(\{b\})=\{a\},\quad \Pi_R^{-1}(R)=X,
\]
and direct computation yields $\mu(\Pi_R^{-1}(B))=\mu(B)$ for each $B\subseteq R$; in particular
\[
\mu(\Pi_R^{-1}(\{a\}))=\mu(\{b,c\})=1+0=1=\mu(\{a\}).
\]
Hence Axiom~III(b) holds.

For the coupling law with $\eta=0$, since $G=\Delta_X$, for any $B\subseteq X$
\[
\mu^{\otimes2}((B\times X)\cap G)=\sum_{x\in B} m(x)^2=\sum_{x\in B} m(x)=\mu(B),
\]
using $m(x)\in\{0,1\}$ so $m(x)^2=m(x)$. Hence Axiom~III(c) holds. All of Axiom~III is therefore satisfied, while Axiom~I fails.
\end{proof}

\begin{proposition}[Independence of Axiom II]\label{prop:indep-II}
There exists a pre-structural datum $\mathcal D_{\neg \mathrm{II}}$ over the ambient signature of Definition~\ref{def:structural-datum} which satisfies Axioms~I and~III but fails Axiom~II.
\end{proposition}

\begin{proof}
Let $X=\{a,b,c\}$ and $\mathcal A=\mathcal P(X)$. Set $R:=X$, $I:=\varnothing$, and $\Pi_R:=\mathrm{id}_X$. Then $\Pi_R$ is idempotent, $\Pi_R|_R=\mathrm{id}_R$, and preimages of measurable subsets are measurable. Hence Axiom~I holds.

Define
\[
G:=\{(a,a),(b,b),(c,c),(a,b),(b,a),(b,c),(c,b)\}\subseteq X\times X.
\]
Then $G$ is reflexive ($\Delta_X\subseteq G$) and symmetric. However, $G$ is not idempotent: the pair $(a,c)$ is obtained via
\[
(a,b)\in G,\quad (b,c)\in G \;\Longrightarrow\; (a,c)\in G\circ G,
\]
but $(a,c)\notin G$, so $G\circ G\supsetneq G$. Hence Axiom~II fails.

Define $\mu$ by $m(a):=1$, $m(b):=0$, $m(c):=1$, and let $\mu^{\otimes2}$ be the atomic product charge. Set $\eta:=0$ and $E_0:=\mu(R)+\mu(I)=2>0$, so Axiom~III(a) holds. Since $\Pi_R=\mathrm{id}_X$, projection--measure invariance (Axiom~III(b)) is trivial.

For the coupling law, we compute $\mu^{\otimes2}((B\times X)\cap G)$ for each singleton $B$. The nonzero contributions come only from pairs $(x,y)\in G$ with $m(x)m(y)\ne 0$, i.e.\ with $x,y\in\{a,c\}$. The only such pairs in $G$ are $(a,a)$ and $(c,c)$ (since $(a,b),(b,c)$ etc.\ involve $b$). Hence for each $B\subseteq X$:
\begin{align*}
B=\{a\}&:\ \mu^{\otimes2}((B\times X)\cap G)=m(a)^2=1=\mu(B),\\
B=\{b\}&:\ \mu^{\otimes2}((B\times X)\cap G)=0=\mu(B),\\
B=\{c\}&:\ \mu^{\otimes2}((B\times X)\cap G)=m(c)^2=1=\mu(B),
\end{align*}
and finite additivity extends this to all $B\subseteq X$. With $\eta=0$ and $\Pi_R=\mathrm{id}_X$, the coupling law reduces to $\mu^{\otimes2}((B\times X)\cap G)=\mu(B)$, which we have verified. Hence Axiom~III(c) holds. All of Axiom~III is therefore satisfied, while Axiom~II fails.
\end{proof}

\begin{remark}[On the $\mu^{\otimes2}$-invisibility of the failure of Axiom~II]\label{rem:II-invisibility}
In the datum $\mathcal D_{\neg \mathrm{II}}$ above, the failure of $G\circ G=G$ occurs at the pair $(a,c)$, which carries positive $\mu^{\otimes2}$-weight ($m(a)m(c)=1$) but lies outside $G$. The separating witness is set-theoretic ($G\circ G\supsetneq G$ as subsets of $X\times X$), not measure-theoretic. This reflects the fact that Axiom~II is a structural rather than measure-theoretic statement; the failure is visible at the level of the relation $G$ itself, regardless of whether the measure detects it through the coupling law.
\end{remark}

\begin{proposition}[Independence of Axiom III]\label{prop:indep-III}
There exists a pre-structural datum $\mathcal D_{\neg \mathrm{III}}$ over the ambient signature of Definition~\ref{def:structural-datum} which satisfies Axioms~I and~II but fails Axiom~III.
\end{proposition}

\begin{proof}
Let $X=\{r,i\}$ and $\mathcal A=\mathcal P(X)$. Set $R:=\{r\}$, $I:=\{i\}$, and define $\Pi_R:X\to R$ by $\Pi_R(r):=r$, $\Pi_R(i):=r$. Then $\Pi_R$ is idempotent, $\Pi_R|_R=\mathrm{id}_R$, and preimages of measurable subsets of $R$ are measurable. Hence Axiom~I holds.

Set $G:=\Delta_X$, satisfying Axiom~II.

Define $\mu$ by $m(r):=m(i):=1$, and let $\mu^{\otimes2}$ be the atomic product charge. Set $\eta:=0$ and $E_0:=\mu(R)+\mu(I)=2>0$, so Axiom~III(a) (conservation) holds.

However, Axiom~III(b) (projection--measure invariance) fails: taking $B=\{r\}\subseteq R$,
\[
\Pi_R^{-1}(\{r\})=X=\{r,i\},
\qquad
\mu(\Pi_R^{-1}(\{r\}))=m(r)+m(i)=2,
\]
while $\mu(\{r\})=1$. Hence $\mu(\Pi_R^{-1}(B))\ne\mu(B)$, and Axiom~III(b) is violated. Consequently Axiom~III as a whole fails.
\end{proof}

\begin{remark}[$\mathcal D_{\neg \mathrm{III}}$ as the reason Axiom~III(b) is needed in the reduction]\label{rem:indep-III-as-counterexample}
The datum $\mathcal D_{\neg \mathrm{III}}$ of Proposition~\ref{prop:indep-III} plays a second role beyond separating Axiom~III: it witnesses the necessity of the Axiom~III(b) hypothesis in part~(iii) of Theorem~\ref{thm:quotient-factorization}. Explicitly, $\mathcal D_{\neg \mathrm{III}}$ satisfies Axioms~I, II, III(a), III(c), fiber measurability, the regularity hypothesis (R--fin), and the full partition $X=R\sqcup I$; yet it is not admissible (because III(b) fails), while its restriction $\mathcal D_{\neg\mathrm{III}}|_R$ is the one--point identity--retraction datum $(\{r\},\{r\})$ with total mass $1$ and trivial relation $\Delta_{\{r\}}$, which is admissible. Hence without the Axiom~III(b) hypothesis, admissibility of $\mathcal D|_R$ would not imply admissibility of $\mathcal D$.
\end{remark}

\subsection{Internal independence of Axiom III}\label{sec:indep-III-internal}

Axiom~III comprises three distinct assertions:
\begin{itemize}[nosep]
\item[(a)] \emph{conservation}: $\mu(R)+\mu(I)=E_0>0$;
\item[(b)] \emph{projection--measure invariance}: $\mu(\Pi_R^{-1}(B))=\mu(B)$ for all measurable $B\subseteq R$;
\item[(c)] \emph{coupling law}: $\mu^{\otimes2}((B\times X)\cap G)=\mu(B)+\eta\,\mu^{\otimes2}((\Pi_R^{-1}(B)\times X)\cap G)$ for all $B\in\mathcal A$.
\end{itemize}

We show that these three subclauses are independent: no one of them follows from Axioms~I, II together with the other two.

\begin{proposition}[Internal independence of Axiom III]\label{prop:indep-III-internal}
Subclauses {\rm(a)}, {\rm(b)}, {\rm(c)} of Axiom~III are mutually independent, relative to Axioms~I and~II.
\end{proposition}

\begin{proof}
We exhibit three structures satisfying Axioms~I,~II and exactly two of {\rm(a), (b), (c)}.

\emph{(a) is independent.} Let $X=\{r\}$, $R=\{r\}$, $I=\varnothing$, $\Pi_R=\mathrm{id}_X$, $G=\Delta_X$, and $\mu\equiv 0$ (the zero charge), with $\eta=0$. Axioms~I and~II hold. Subclause~(b) is trivial ($\Pi_R=\mathrm{id}$), and subclause~(c) holds (both sides equal $0$). However, $\mu(R)+\mu(I)=0$, so no $E_0>0$ satisfies conservation; subclause~(a) fails. This shows (a) is not derivable from Axioms~I, II, (b), (c).

\emph{(b) is independent.} The structure $\mathcal D_{\neg \mathrm{III}}$ of Proposition~\ref{prop:indep-III} satisfies Axioms~I,~II and subclause~(a) (with $E_0=2$), and a direct check (as in the proof of Proposition~\ref{prop:indep-III}) shows that the coupling law~(c) holds with $\eta=0$: for $B=\{r\}$, $\mu^{\otimes2}((\{r\}\times X)\cap \Delta_X)=m(r)^2=1=\mu(\{r\})$, and similarly for $B=\{i\}$ and $B=X$. However, subclause~(b) fails as shown in Proposition~\ref{prop:indep-III}. Hence (b) is not derivable from Axioms~I,~II,~(a),~(c).

\emph{(c) is independent.} Let $X=\{a,b\}$, $R=X$, $I=\varnothing$, $\Pi_R=\mathrm{id}_X$, $G=X\times X$, with point weights $m(a)=1$, $m(b)=0$, and $\eta=1/2$. Axioms~I and~II hold (the total relation $X\times X$ is reflexive, symmetric, and idempotent). Subclause~(a): $\mu(R)+\mu(I)=1>0$, take $E_0=1$. Subclause~(b): $\Pi_R=\mathrm{id}_X$ makes invariance trivial. However, for $B=\{a\}$,
\[
\mu^{\otimes2}((\{a\}\times X)\cap (X\times X))=m(a)\bigl(m(a)+m(b)\bigr)=1\cdot 1=1,
\]
while
\[
\mu(B)+\eta\,\mu^{\otimes2}((\Pi_R^{-1}(B)\times X)\cap G)=1+\tfrac{1}{2}\cdot 1=\tfrac{3}{2}\ne 1,
\]
so the coupling law fails. Hence (c) is not derivable from Axioms~I,~II,~(a),~(b).
\end{proof}

\subsection{The main independence theorem and minimality corollary}\label{sec:indep-main}

\begin{theorem}[Independence]\label{thm:independence-main}
The axioms $\mathrm{I}$, $\mathrm{II}$, $\mathrm{III}$ of Section~\ref{sec:axioms} are mutually independent: for each $\mathsf X\in\{\mathrm{I},\mathrm{II},\mathrm{III}\}$ there exists a pre-structural datum over the ambient signature of Definition~\ref{def:structural-datum} satisfying all axioms except $\mathsf X$. Moreover, the three subclauses {\rm(a), (b), (c)} of Axiom~III are themselves mutually independent, relative to Axioms~I and~II.
\end{theorem}

\begin{proof}
Combine Propositions~\ref{prop:indep-I}, \ref{prop:indep-II}, \ref{prop:indep-III}, and~\ref{prop:indep-III-internal}.
\end{proof}

\begin{corollary}[Minimality]\label{cor:minimality}
The axiom system $\{\mathrm{I},\mathrm{II},\mathrm{III}\}$ is minimal in the following sense: removing any one axiom, or any one subclause of Axiom~III, yields a strictly weaker system admitting structures not admitted by the full system.
\end{corollary}

\begin{proof}
Immediate from Theorem~\ref{thm:independence-main}: the separating models of Propositions~\ref{prop:indep-I}, \ref{prop:indep-II}, \ref{prop:indep-III}, \ref{prop:indep-III-internal} witness the claim in each case.
\end{proof}

\begin{remark}[On the baseline analytic structure]\label{rem:baseline-analytic}
The independence results above are stated relative to the admitted analytic background: a set algebra $\mathcal A\subseteq\mathcal P(X)$, a finitely additive measure $\mu$, a product charge $\mu^{\otimes2}$ satisfying the rectangle rule, and a retraction $\Pi_R:X\to R$ together with a binary relation $G\subseteq X\times X$. We do not claim purely logical independence stripped of this analytic setting; such a claim would require encoding the axioms in a formal system without those primitives, which is beyond the scope of this paper.
\end{remark}

\begin{remark}[Roles of the three axioms]\label{rem:roles}
The three separating models clarify the distinct structural roles of the three axioms:
\begin{itemize}[nosep]
\item Axiom~I governs the \emph{projective skeleton} of the datum (retraction data and measurability of preimages). Its failure in $\mathcal D_{\neg\mathrm{I}}$ is visible already at the level of $\Pi_R$, before measure-theoretic content is invoked.
\item Axiom~II governs the \emph{closure structure} of the interaction relation. Its failure in $\mathcal D_{\neg\mathrm{II}}$ is set-theoretic ($G\circ G\supsetneq G$), independent of measure.
\item Axiom~III governs the \emph{quantitative coupling} between measure, projection, and relation. Its three subclauses separate cleanly: conservation fixes total mass, invariance governs pointwise behavior under $\Pi_R$, and the coupling law governs the two-point curvature load.
\end{itemize}
\end{remark}

% -------------------------------------------------------

\section{\texorpdfstring{$\sigma$}{sigma}--additive refinement and constraints}\label{sec:sigma-additive}

This paper is formulated for finitely additive measures in order to keep the axioms minimal and
to ensure broad model existence. It is natural to ask what changes if one
requires $\mu$ to be $\sigma$--additive (and hence admits the standard product measure by the classical construction).

\subsection{Setup}
Assume $(X,\Sigma,\mu)$ is a $\sigma$--additive measure space (see e.g.\ \cite{Halmos,Bogachev}). In order to speak unambiguously about the product measure
$\mu^{\otimes2}$ on $(X\times X,\Sigma\otimes\Sigma)$ we assume $\mu$ is $\sigma$--finite when needed.
In that case the standard product measure exists and is uniquely characterized by
\[
\mu^{\otimes2}(A\times B)=\mu(A)\mu(B)\qquad(A,B\in\Sigma),
\]
together with $\sigma$--additivity on $\Sigma\otimes\Sigma$.

We emphasize that the conclusions below use only the $B=X$ instance of the coupling law,
and thus apply in any $\sigma$--additive setting where $\mu^{\otimes2}$ is defined.

\subsection{A universal constraint from the coupling law}
The coupling law in Axiom~III implies a strong global constraint already at $B=X$.

\begin{proposition}[Global constraint at $B=X$]\label{prop:global-constraint}
Assume $\mu$ is $\sigma$--additive and $\mu^{\otimes2}$ is the standard product measure.
If the coupling law holds for $B=X$, then
\[
(1-\eta)\,\mu^{\otimes2}(G)=\mu(X).
\]
In particular, if $\mu(X)\in(0,\infty)$ then
\[
\mu^{\otimes2}(G)=\frac{\mu(X)}{1-\eta}.
\]
\end{proposition}

\begin{proof}
Take $B=X$ in the coupling law:
\[
\mu^{\otimes2}(G)=\mu(X)+\eta\,\mu^{\otimes2}((\Pi_R^{-1}(X)\times X)\cap G).
\]
Since $(X\times X)\cap G=G$ and $\Pi_R^{-1}(X)=X$, this becomes
$\mu^{\otimes2}(G)=\mu(X)+\eta\,\mu^{\otimes2}(G)$, i.e. $(1-\eta)\mu^{\otimes2}(G)=\mu(X)$.
\end{proof}

\begin{corollary}[Finite measure bound]\label{cor:finite-measure-bound}
Assume $\mu(X)=M\in(0,\infty)$. Then the coupling law at $B=X$ forces
\[
\frac{1}{1-\eta}\le M.
\]
Equivalently, $\eta\le 1-\frac{1}{M}$.
\end{corollary}

\begin{proof}
By Proposition~\ref{prop:global-constraint}, $\mu^{\otimes2}(G)=M/(1-\eta)$.
Since $\mu^{\otimes2}(G)\le \mu^{\otimes2}(X\times X)=M^2$, we obtain $M/(1-\eta)\le M^2$, i.e. $1/(1-\eta)\le M$.
\end{proof}

\begin{corollary}[Probability case forces $\eta=0$]\label{cor:probability-forces-eta0}
If $\mu$ is a probability measure, i.e. $\mu(X)=1$, then the coupling law at $B=X$ implies
\[
\eta=0
\quad\text{and}\quad
\mu^{\otimes2}(G)=1,
\]
so $G$ has full $\mu^{\otimes2}$--measure in $X\times X$.
\end{corollary}

\begin{proof}
From Proposition~\ref{prop:global-constraint}, $\mu^{\otimes2}(G)=1/(1-\eta)\ge 1$.
But $\mu^{\otimes2}(G)\le \mu^{\otimes2}(X\times X)=1$, hence $\mu^{\otimes2}(G)=1$ and $\eta=0$.
\end{proof}

\begin{proposition}[Feasibility window for $\eta\neq 0$ under finite total mass]\label{prop:eta-window}
Assume $\mu$ is $\sigma$--additive, $\sigma$--finite, and $\mu(X)=M\in(0,\infty)$.
If the coupling law holds at $B=X$, then necessarily
\[
0\le \eta \le 1-\frac{1}{M},
\]
and moreover $\mu^{\otimes2}(G)=\frac{M}{1-\eta}\in(0,M^2]$.
In particular, $\eta\neq 0$ is possible only when $M>1$ (i.e.\ the measure is not probability-normalized).
\end{proposition}

\begin{proof}
The bound on $\eta$ is Corollary~\ref{cor:finite-measure-bound}.
The identity for $\mu^{\otimes2}(G)$ follows from Proposition~\ref{prop:global-constraint}.
\end{proof}

\begin{remark}[Infinite-mass regimes]\label{rem:infinite-mass}
If $\mu(X)=+\infty$, Proposition~\ref{prop:global-constraint} becomes $+\infty=(1-\eta)\mu^{\otimes2}(G)$, which does not by itself
exclude $\eta\neq 0$. In such regimes one must use additional instances of the coupling law (i.e.\ varying $B$) or add normalization
assumptions to obtain effective constraints.
\end{remark}

\begin{remark}[Scaling intuition]\label{rem:eta-scaling}
\sloppy
In an approximating sequence of finite models $(X_n,\allowbreak \mu_n,\allowbreak G_n,\allowbreak \Pi_{R,n},\allowbreak \eta_n)$, the limiting behavior of $\eta_n$ is fundamentally normalization--dependent. Under standard probability normalization ($\mu_n(X_n)=1$), the global constraint typically forces $\eta_n\to 0$, often rendering the coupling trivial. Conversely, in regimes where $\mu_n(X_n)=M_n\to M>1$, $\eta_n$ can stabilize to a nonzero constant compatible with the total mass $M$. In large--mass regimes ($M_n\to\infty$), the global identity does not by itself force $\eta_n\to 0$, shifting the structural burden to local instances of the coupling law (varying $B$), which become essential for characterizing the infinitesimal structure of the limit.
\end{remark}

\paragraph{Further directions (minimal).}
(i) Study additional instances of the coupling law (varying $B$) to obtain nontrivial constraints beyond the $B=X$ global identity.
(ii) Classify $\eta\neq 0$ solutions under normalization choices (finite $M$, infinite mass, or $\sigma$--finite restrictions) and determine when non-diagonal $G$ must have full or intermediate $\mu^{\otimes2}$--measure.
(iii) Compare the finitely additive model families of Section~5 with $\sigma$--additive realizability, identifying the exact obstruction created by probability normalization.

\begin{remark}[Interpretation]
Corollary~\ref{cor:probability-forces-eta0} explains why nontrivial $\eta\neq 0$ families are naturally formulated
either (i) in the finitely additive setting, or (ii) for $\sigma$--additive measures that are not normalized as probabilities.
The global constraint arises solely from the $B=X$ instance of the coupling law; under a fixed normalization $\mu(X)=M$, the admissible range of $\eta$ shrinks as $M$ decreases, and any sequence of models in which $\eta$ is to stabilize away from zero therefore requires $M$ to grow accordingly.
\end{remark}

% -------------------------------------------------------

\section{Categorical structure}\label{sec:category}

We record a basic categorical organization of structural models: the appropriate notion of morphism, the resulting category, and the preservation of the coupling law under morphisms.

\subsection{\texorpdfstring{Morphisms and the category $\mathbf{Struct}$}{Morphisms and the category Struct}}

Throughout this section we fix the scalar $\eta\in[0,1)$, but treat $E_0\in(0,\infty)$ as \emph{object data}: it is recorded as part of each structural model and is allowed to vary between objects. Morphisms interact with $E_0$ via~(M4) below, which yields an automatic inequality $E_0'\le E_0$ (see Remark~\ref{rem:E0-morphism}), rather than an a priori identification. From this point on, $\mathbf{Struct}_{\eta}$ denotes the category whose objects are admissible structural models (Definition~\ref{def:admissible}) with common $\eta\in[0,1)$, and whose morphisms are defined below.

\begin{definition}[Morphism of structural models]\label{def:morphism}
Let $\mathcal M=(X,\allowbreak\mathcal A,\allowbreak\mu,\allowbreak\mu^{\otimes2},\allowbreak R,\allowbreak I,\allowbreak\Pi_R,\allowbreak G,\allowbreak E_0,\allowbreak\eta)$ and $\mathcal M'=(X',\allowbreak\mathcal A',\allowbreak\mu',\allowbreak\mu'^{\otimes2},\allowbreak R',\allowbreak I',\allowbreak\Pi_{R'},\allowbreak G',\allowbreak E_0',\allowbreak\eta)$ be admissible structural models (Definition~\ref{def:admissible}) sharing a common value of $\eta\in[0,1)$; the scalars $E_0,E_0'$ are part of the object data and need not agree a priori. A \emph{morphism} $\phi:\mathcal M\to\mathcal M'$ is a map $\phi:X\to X'$ satisfying:
\begin{enumerate}[label=\emph{(M\arabic*)},nosep,leftmargin=2.5em]
\item \emph{Measurability}: $\phi^{-1}(B')\in\mathcal A$ for every $B'\in\mathcal A'$, and $(\phi\times\phi)^{-1}(S')\in\mathcal A\otimes\mathcal A$ for every $S'\in\mathcal A'\otimes\mathcal A'$.
\item \emph{Projection commutativity}: $\phi\circ\Pi_R=\Pi_{R'}\circ\phi$.
\item \emph{Relation preservation}: $(\phi\times\phi)(G)\subseteq G'$, i.e., $(x,y)\in G$ implies $(\phi(x),\phi(y))\in G'$.
\item \emph{Measure--preservation}: $\mu(\phi^{-1}(B'))=\mu'(B')$ for every $B'\in\mathcal A'$, and $\mu^{\otimes2}((\phi\times\phi)^{-1}(S'))=\mu'^{\otimes2}(S')$ for every $S'\in\mathcal A'\otimes\mathcal A'$.
\end{enumerate}
\end{definition}

\begin{remark}[$E_0$ under morphisms]\label{rem:E0-morphism}
Axiom~(M4) alone controls the relation between source and target scalars:
\[
E_0'=\mu'(R')+\mu'(I')=\mu(\phi^{-1}(R'))+\mu(\phi^{-1}(I'))=\mu\big(\phi^{-1}(R'\cup I')\big)\le\mu(X).
\]
If $\mathcal M$ has full partition ($X=R\sqcup I$), then $\mu(X)=E_0$ and hence $E_0'\le E_0$. Equality $E_0'=E_0$ holds when, in addition, $\phi^{-1}(R'\cup I')=X$, for instance when $\phi$ is surjective and $\mathcal M'$ also has full partition. Morphisms are thus in general mass--non--increasing, and mass--preserving in the natural full--partition surjective case. No separate scalar--identification axiom is imposed, since the relation between $E_0$ and $E_0'$ is already controlled by~(M4).
\end{remark}

\begin{remark}[Automatic consequences]\label{rem:morphism-consequences}
Three consequences follow directly from the definition:
\begin{itemize}[nosep]
\item $\phi(R)\subseteq R'$, since $\phi(R)=\phi(\Pi_R(X))=\Pi_{R'}(\phi(X))\subseteq R'$ by~(M2) and Axiom~I for $\mathcal M'$.
\item For $B'\subseteq R'$ measurable, $\mu'(B')=\mu(\phi^{-1}(B'))$, which combined with projection--measure invariance in $\mathcal M'$ recovers invariance in $\mathcal M$ along $\phi$.
\item The pushforward $\phi_*\mu^{\otimes2}$ coincides with $\mu'^{\otimes2}$ on rectangles, since
\[
(\phi_*\mu^{\otimes2})(A'\times B')=\mu^{\otimes2}\bigl(\phi^{-1}(A')\times\phi^{-1}(B')\bigr)=\mu(\phi^{-1}(A'))\,\mu(\phi^{-1}(B'))
\]
which by~(M4) equals $\mu'(A')\,\mu'(B')=\mu'^{\otimes2}(A'\times B')$, using the rectangle rule.
\end{itemize}
\end{remark}

\begin{proposition}[Composition and identity]\label{prop:category}
Morphisms of structural models compose: if $\phi:\mathcal M\to\mathcal M'$ and $\psi:\mathcal M'\to\mathcal M''$ are morphisms (sharing a common $\eta$), then $\psi\circ\phi:\mathcal M\to\mathcal M''$ is a morphism. Moreover, the identity map $\mathrm{id}_X:\mathcal M\to\mathcal M$ is a morphism. Hence structural models with a common value of $\eta\in[0,1)$ (with $E_0$ allowed to vary as object data) together with their morphisms form a category, which we denote by $\mathbf{Struct}_{\eta}$.
\end{proposition}

\begin{proof}
We check the four morphism axioms for $\psi\circ\phi$.

(M1) Measurability of $\psi\circ\phi$ follows from $(\psi\circ\phi)^{-1}(B'')=\phi^{-1}(\psi^{-1}(B''))\in\mathcal A$ since $\psi^{-1}(B'')\in\mathcal A'$ and $\phi^{-1}(\mathcal A')\subseteq\mathcal A$. The product--measurability condition follows by the same argument applied to $(\psi\times\psi)\circ(\phi\times\phi)$.

(M2) Projection commutativity:
\[
(\psi\circ\phi)\circ\Pi_R=\psi\circ(\phi\circ\Pi_R)=\psi\circ\Pi_{R'}\circ\phi=\Pi_{R''}\circ\psi\circ\phi=\Pi_{R''}\circ(\psi\circ\phi).
\]

(M3) Relation preservation:
\[
((\psi\circ\phi)\times(\psi\circ\phi))(G)=(\psi\times\psi)((\phi\times\phi)(G))\subseteq(\psi\times\psi)(G')\subseteq G''.
\]

(M4) Measure--preservation, for $B''\in\mathcal A''$:
\[
\mu((\psi\circ\phi)^{-1}(B''))=\mu(\phi^{-1}(\psi^{-1}(B'')))=\mu'(\psi^{-1}(B''))=\mu''(B'').
\]
The product--measure identity follows analogously. Identity morphism is trivial.
\end{proof}

\subsection{Preservation of the coupling law}

Morphisms of structural models preserve not only the listed data but also the coupling law derived from Axiom~III.

\begin{proposition}[Coupling law is transported along exact pullback morphisms]\label{prop:coupling-preserved}
Let $\phi:\mathcal M\to\mathcal M'$ be a morphism. Suppose $\mathcal M$ satisfies Axiom~III (coupling law). Then:
\begin{enumerate}[label=\emph{(\roman*)},nosep]
\item \emph{(Source--side inequality)} Under the hypothesis $(\phi\times\phi)^{-1}(G')\subseteq G$, for every $B'\in\mathcal A'$,
\[
\mu'^{\otimes2}\bigl((B'\times X')\cap G'\bigr)\ \le\ \mu^{\otimes2}\bigl((\phi^{-1}(B')\times X)\cap G\bigr),
\]
and the right--hand side admits the coupling--law expansion for $\mathcal M$:
\[
\mu^{\otimes2}\bigl((\phi^{-1}(B')\times X)\cap G\bigr)=\mu(\phi^{-1}(B'))+\eta\,\mu^{\otimes2}\bigl((\Pi_R^{-1}(\phi^{-1}(B'))\times X)\cap G\bigr).
\]
\item \emph{(Equality under exact pullback and surjectivity)} If in addition $\phi$ is surjective and $(\phi\times\phi)^{-1}(G')=G$, then $\mathcal M'$ satisfies the coupling law in full:
\[
\mu'^{\otimes2}\bigl((B'\times X')\cap G'\bigr)=\mu'(B')+\eta\,\mu'^{\otimes2}\bigl((\Pi_{R'}^{-1}(B')\times X')\cap G'\bigr)\qquad(\forall B'\in\mathcal A').
\]
\end{enumerate}
\end{proposition}

\begin{proof}
Let $B'\in\mathcal A'$. By~(M4),
\[
\mu'^{\otimes2}((B'\times X')\cap G')=\mu^{\otimes2}\bigl((\phi\times\phi)^{-1}((B'\times X')\cap G')\bigr).
\]
Since $(\phi\times\phi)^{-1}(B'\times X')=\phi^{-1}(B')\times X$, we have
\[
(\phi\times\phi)^{-1}((B'\times X')\cap G')=(\phi^{-1}(B')\times X)\cap(\phi\times\phi)^{-1}(G').
\]
\emph{Part (i).} Under $(\phi\times\phi)^{-1}(G')\subseteq G$, the set on the right is contained in $(\phi^{-1}(B')\times X)\cap G$; by monotonicity of $\mu^{\otimes2}$,
\[
\mu'^{\otimes2}((B'\times X')\cap G')\le\mu^{\otimes2}((\phi^{-1}(B')\times X)\cap G).
\]
By the coupling law for $\mathcal M$ applied to $B=\phi^{-1}(B')$,
\[
\mu^{\otimes2}((\phi^{-1}(B')\times X)\cap G)=\mu(\phi^{-1}(B'))+\eta\,\mu^{\otimes2}((\Pi_R^{-1}(\phi^{-1}(B'))\times X)\cap G),
\]
giving the inequality asserted in (i), with the right--hand side admitting the stated expansion.

\emph{Part (ii).} Under the stronger hypotheses $(\phi\times\phi)^{-1}(G')=G$ and $\phi$ surjective, the inclusion in the displayed identity becomes equality:
\[
(\phi\times\phi)^{-1}((B'\times X')\cap G')=(\phi^{-1}(B')\times X)\cap G,
\]
and hence
\[
\mu'^{\otimes2}((B'\times X')\cap G')=\mu^{\otimes2}((\phi^{-1}(B')\times X)\cap G).
\]
Applying the coupling law for $\mathcal M$ to $B=\phi^{-1}(B')$, using (M2) to rewrite
\[
\Pi_R^{-1}(\phi^{-1}(B'))=\phi^{-1}(\Pi_{R'}^{-1}(B')),
\]
and applying (M4) for both $\mu$ and $\mu^{\otimes2}$:
\begin{align*}
\mu'^{\otimes2}((B'\times X')\cap G')
&=\mu(\phi^{-1}(B'))+\eta\,\mu^{\otimes2}((\phi^{-1}(\Pi_{R'}^{-1}(B'))\times X)\cap G)\\
&=\mu'(B')+\eta\,\mu'^{\otimes2}((\Pi_{R'}^{-1}(B')\times X')\cap G'),
\end{align*}
which is the coupling law in $\mathcal M'$, valid for every $B'\in\mathcal A'$ (using surjectivity of $\phi$ to range over all $B'$).
\end{proof}

\begin{remark}[On the hypotheses]\label{rem:surjectivity}
Part~(ii) of Proposition~\ref{prop:coupling-preserved} is the intended ``transport of coupling law'' statement; the strong hypothesis $(\phi\times\phi)^{-1}(G')=G$ combined with surjectivity of $\phi$ is automatic, for instance, when $\mathcal M$ is the full pullback of $\mathcal M'$ along $\phi$ (i.e., when $\mathcal M$ is obtained from $\mathcal M'$ by pulling back the measure and relation along $\phi$). In the weaker inclusion case of part~(i), only an inequality survives, and the coupling law of $\mathcal M'$ is bounded by the coupling--law expression of $\mathcal M$ on the pulled--back test set. The discrepancy between (i) and (ii) records precisely the information lost when $\phi$ fails to be surjective or the relation is preserved only inclusively.
\end{remark}

\subsection{\texorpdfstring{Idempotent endomorphisms in $\mathbf{Struct}$}{Idempotent endomorphisms in Struct}}

The axioms admit a natural categorical reformulation in terms of idempotents.

\begin{proposition}[Retraction as a split idempotent on objects]\label{prop:retraction-idempotent}
Let $\mathcal M=(X,\allowbreak\mathcal A,\allowbreak\mu,\allowbreak\mu^{\otimes2},\allowbreak R,\allowbreak I,\allowbreak\Pi_R,\allowbreak G,\allowbreak E_0,\allowbreak\eta)$ be a structural model. The map $\Pi_R:X\to X$ is idempotent ($\Pi_R\circ\Pi_R=\Pi_R$) at the level of the underlying measurable-relation data. Restricting codomain to $R\in\mathcal A$, $\Pi_R$ factors as
\[
\Pi_R\;=\;\iota_R\circ\rho_R,
\]
where $\rho_R:X\to R$ is the surjection defined by $\Pi_R$ and $\iota_R:R\hookrightarrow X$ is the inclusion. The equality $\rho_R\circ\iota_R=\mathrm{id}_R$ holds by Axiom~I.
\end{proposition}

\begin{proof}
Idempotence is Axiom~I. The factorization $\Pi_R=\iota_R\circ\rho_R$ is definitional: $\rho_R$ is simply $\Pi_R$ with codomain restricted to $R$, and $\iota_R:R\to X$ is the set--theoretic inclusion. The identity $\rho_R\circ\iota_R=\mathrm{id}_R$ follows from $\Pi_R|_R=\mathrm{id}_R$ (Axiom~I).
\end{proof}

\begin{proposition}[Relation idempotence at the morphism level]\label{prop:G-idempotent}
The relation $G\subseteq X\times X$, viewed as a binary endorelation, is an idempotent endomorphism in the poset category of binary relations on $X$ under relational composition: $G\circ G=G$ (Axiom~II). Consequently $G$ defines an equivalence relation $\sim_G$ on $X$.
\end{proposition}

\begin{proof}
Idempotence is Axiom~II; Proposition~\ref{prop:elementary}(iii) upgrades reflexive, symmetric, idempotent relations to equivalence relations.
\end{proof}

\begin{remark}[Monad--comonad perspective]\label{rem:monad-comonad}
Propositions~\ref{prop:retraction-idempotent} and~\ref{prop:G-idempotent} together place the two structural idempotents of a structural model in analogy with the two sides of an idempotent monad--comonad pair:
\begin{itemize}[nosep]
\item The retraction $\Pi_R$ plays the role of the monad component: it is an idempotent endomorphism with a distinguished inclusion section $\iota_R$.
\item The relation $G$ plays the role of the comonad component: its idempotence generates equivalence classes that organize $X$ into a quotient structure.
\end{itemize}
The coupling law of Axiom~III then expresses a quantitative compatibility between these two idempotent structures, mediated by the measure $\mu$ and the parameter $\eta$.
\end{remark}

\subsection{\texorpdfstring{The subcategory of $\Pi_R=\mathrm{id}$ models}{The subcategory of Pi\_R=id models}}

Within $\mathbf{Struct}_{\eta}$, the condition $\Pi_R=\mathrm{id}_X$ singles out a structurally simpler subcategory.

\begin{definition}[Subcategory of identity--retraction models]\label{def:struct-id}
Let $\mathbf{Struct}^{\mathrm{id}}_{\eta}\subseteq\mathbf{Struct}_{\eta}$ be the full subcategory consisting of models $\mathcal M$ with $\Pi_R=\mathrm{id}_X$ (so $R=X$, $I=\varnothing$).
\end{definition}

\begin{proposition}[Restriction functor, under fiber hypotheses]\label{prop:restriction-functor}
Let $\mathbf{Struct}_{\eta}^{\mathrm{fib}}\subseteq\mathbf{Struct}_{\eta}$ denote the full subcategory whose objects are those structural models $\mathcal M$ satisfying the measurability hypothesis of Proposition~\ref{prop:fiber-annihilation} (namely $\{r\}\in\mathcal A$ and $F_r=\Pi_R^{-1}(\{r\})\in\mathcal A$ for every $r\in R$) together with one of the regularity hypotheses \emph{(R--fin)} or \emph{(R--ctbl)} of that proposition. Then the assignment
\[
\rho:\mathbf{Struct}_{\eta}^{\mathrm{fib}}\longrightarrow\mathbf{Struct}^{\mathrm{id}}_{\eta},\qquad
\mathcal M\mapsto\mathcal M|_R,\quad \phi\mapsto\phi|_R,
\]
defines a functor, where $\mathcal M|_R$ is the model with underlying set $R$, set algebra $\mathcal A\cap\mathcal P(R)$, measure $\mu|_R$, product charge $\mu^{\otimes2}|_{R\times R}$, trivial retraction $\Pi_{R|_R}=\mathrm{id}_R$, restricted relation $G|_R:=G\cap(R\times R)$, scalar $E_0(\mathcal M|_R):=\mu(R)$, and coupling parameter $\eta$.
\end{proposition}

\begin{proof}
We verify that $\mathcal M|_R$ satisfies Axioms~I--III, so that $\mathcal M|_R\in\mathbf{Struct}^{\mathrm{id}}_{\eta}$. The argument relies only on Proposition~\ref{prop:fiber-annihilation} (the fiber--annihilation consequence $\mu(X\setminus R)=0$ under the hypotheses defining $\mathbf{Struct}_{\eta}^{\mathrm{fib}}$) and on Axioms~I--III for $\mathcal M$; it does not invoke Theorem~\ref{thm:quotient-factorization}.

\emph{Axiom I on $\mathcal M|_R$.} Immediate from $\Pi_{R|_R}=\mathrm{id}_R$: the retraction is idempotent, and $\Pi_{R|_R}|_R=\mathrm{id}_R$ trivially, with measurability of preimages holding on $\mathcal A\cap\mathcal P(R)$.

\emph{Axiom II on $\mathcal M|_R$.} The restricted relation $G|_R:=G\cap(R\times R)$ is reflexive on $R$ (since $\Delta_R\subseteq\Delta_X\subseteq G$ and $\Delta_R\subseteq R\times R$) and symmetric by restriction. For idempotence, observe first that $G|_R\circ G|_R\subseteq G\circ G\cap(R\times R)=G\cap(R\times R)=G|_R$ (using Axiom~II for $\mathcal M$ and the fact that composing relations on $R\times R$ stays in $R\times R$). Conversely, if $(x,z)\in G|_R$, then $x,z\in R$ and, by reflexivity of $G|_R$ giving $(x,x)\in G|_R$, we have $(x,z)=(x,x)\circ(x,z)\in G|_R\circ G|_R$. Therefore $G|_R\circ G|_R=G|_R$.

\emph{Axiom III(a) on $\mathcal M|_R$.} We must show $\mu(R)>0$. By Proposition~\ref{prop:fiber-annihilation}, $\mu(X\setminus R)=0$. Since $I\subseteq X\setminus R$ by construction, $\mu(I)=0$. Axiom~III(a) for $\mathcal M$ gives $\mu(R)+\mu(I)=E_0(\mathcal M)>0$, hence $\mu(R)=E_0(\mathcal M)>0$. We record this positive scalar as $E_0(\mathcal M|_R):=\mu(R)$.

\emph{Axiom III(b) on $\mathcal M|_R$.} Trivial, since $\Pi_{R|_R}=\mathrm{id}_R$.

\emph{Axiom III(c) on $\mathcal M|_R$.} We show directly that for every $B\in\mathcal A\cap\mathcal P(R)$,
\[
\mu^{\otimes2}((B\times R)\cap G|_R)=\mu(B)+\eta\,\mu^{\otimes2}((B\times R)\cap G|_R),
\]
equivalently $(1-\eta)\mu^{\otimes2}((B\times R)\cap G|_R)=\mu(B)$. Fix $B\in\mathcal A\cap\mathcal P(R)$, so $B\subseteq R$ and hence $\Pi_R^{-1}(B)\cap R=B$ (by $\Pi_R|_R=\mathrm{id}_R$). Applying the coupling law for $\mathcal M$ to this $B$ (viewed as an element of $\mathcal A$):
\[
\mu^{\otimes2}((B\times X)\cap G)=\mu(B)+\eta\,\mu^{\otimes2}((\Pi_R^{-1}(B)\times X)\cap G). \tag{$\ast$}
\]
Since $\mu(X\setminus R)=0$, monotonicity and the rectangle rule give
\[
S\subseteq X\times(X\setminus R)\ \text{measurable}\ \Longrightarrow\ \mu^{\otimes2}(S)\le\mu^{\otimes2}(X\times(X\setminus R))=\mu(X)\,\mu(X\setminus R)=0,
\]
and similarly $\mu^{\otimes2}(S')=0$ for every measurable $S'\subseteq(X\setminus R)\times X$. Consequently,
\[
\mu^{\otimes2}((B\times X)\cap G)=\mu^{\otimes2}((B\times R)\cap G)=\mu^{\otimes2}((B\times R)\cap G|_R),
\]
where the last equality uses $B\subseteq R$ so $(B\times R)\cap G=(B\times R)\cap G\cap(R\times R)=(B\times R)\cap G|_R$. For the right--hand side of~$(\ast)$, note that $\Pi_R^{-1}(B)\setminus B\subseteq X\setminus R$: indeed, if $x\in\Pi_R^{-1}(B)\cap R$ then $\Pi_R(x)=x\in B$ (Axiom~I), so $x\in B$; the contrapositive gives $\Pi_R^{-1}(B)\setminus B\subseteq X\setminus R$. Hence
\[
\mu^{\otimes2}((\Pi_R^{-1}(B)\times X)\cap G)=\mu^{\otimes2}((B\times X)\cap G)=\mu^{\otimes2}((B\times R)\cap G|_R),
\]
using that the ``extra'' set $(\Pi_R^{-1}(B)\setminus B)\times X$ is contained in $(X\setminus R)\times X$, hence $\mu^{\otimes2}$--null by monotonicity as above, and applying the previous reduction. Substituting both sides of~$(\ast)$:
\[
\mu^{\otimes2}((B\times R)\cap G|_R)=\mu(B)+\eta\,\mu^{\otimes2}((B\times R)\cap G|_R),
\]
which is Axiom~III(c) for $\mathcal M|_R$.

\emph{Functoriality.} If $\phi:\mathcal M\to\mathcal M'$ is a morphism in $\mathbf{Struct}_{\eta}^{\mathrm{fib}}$, then $\phi(R)\subseteq R'$ by Remark~\ref{rem:morphism-consequences}, so $\phi|_R:R\to R'$ is well--defined. Axioms~(M1)--(M4) for $\phi|_R$ follow by restriction from those of $\phi$. Composition and identity are preserved, so $\rho$ is a functor.
\end{proof}

\begin{remark}[The restriction functor as a retraction of categories]\label{rem:restriction-retraction}
The subcategory $\mathbf{Struct}^{\mathrm{id}}_{\eta}$ includes canonically into $\mathbf{Struct}_{\eta}^{\mathrm{fib}}$ via the inclusion $\iota$ that regards an identity--retraction object as already satisfying the fiber hypotheses trivially (with $F_r=\{r\}$ for every $r\in R=X$). The composition $\rho\circ\iota$ is the identity on $\mathbf{Struct}^{\mathrm{id}}_{\eta}$: restricting an identity--retraction model to $R=X$ returns the model unchanged. In this precise sense, $\rho$ is a retraction of $\mathbf{Struct}_{\eta}^{\mathrm{fib}}$ onto its subcategory $\mathbf{Struct}^{\mathrm{id}}_{\eta}$, mirroring at the categorical level the object--level retraction $\Pi_R:X\to R$.
\end{remark}

\begin{remark}[Connection to the classification theorem]\label{rem:category-classification}
\sloppy
Within the full subcategory of $\mathbf{Struct}^{\mathrm{id}}_{\eta}$ whose objects satisfy the full hypotheses of Theorem~\ref{thm:pi-id-classification}, that theorem provides a dichotomy: for every $G$--equivalence class $C_k$, either $\mu(C_k)=0$ or $\mu(C_k)=(1-\eta)^{-1}$. Equivalently, every positive--mass class satisfies $\mu(C_k)=(1-\eta)^{-1}$. On that subcategory, each object can be described combinatorially by a partition $(X,(C_k)_{k\in K})$ together with admissible mass data on the classes (namely the dichotomy $\mu(C_k)\in\{0,(1-\eta)^{-1}\}$) and the induced product charge on $X\times X$; morphisms preserve the partition structure and the measure/charge data in the sense of Definition~\ref{def:morphism}. The restriction functor $\rho$ of Proposition~\ref{prop:restriction-functor} extracts this combinatorial core from an arbitrary structural model under the fiber hypotheses, provided the core satisfies those further hypotheses.
\end{remark}

\subsection{Quotient factorization and structural reduction}\label{sec:quotient-factorization}

We now sharpen the restriction functor of Proposition~\ref{prop:restriction-functor} into a structural reduction theorem: under fiber measurability and regularity hypotheses, all curvature--load information relevant to Axioms~I--III reduces to the identity--retraction core $\mathcal M|_R$, and the coupling law on $\mathcal M$ is equivalent to the coupling law on $\mathcal M|_R$. This extends the classification theorem (Theorem~\ref{thm:pi-id-classification}), which addresses the case $\Pi_R=\mathrm{id}_X$, to the general case.

Throughout this subsection let
\[
\mathcal D=(X,\mathcal A,\mu,\mu^{\otimes2},R,I,\Pi_R,G,E_0,\eta)
\]
be a pre-structural datum (Definition~\ref{def:structural-datum}) with $\eta\in[0,1)$ and $\mu(X)<\infty$. Additional assumptions (fiber measurability, regularity, partial axioms, or full admissibility) will be imposed explicitly in each statement.

\begin{definition}[Fiber partition]\label{def:fibers}
For each $r\in R$, the \emph{fiber of $\Pi_R$ over $r$} is
\[
F_r:=\Pi_R^{-1}(\{r\})=\{x\in X:\Pi_R(x)=r\}.
\]
Since $\Pi_R$ is a function into $R$, the fibers $\{F_r\}_{r\in R}$ always partition $X$ as a disjoint union indexed by the image of $\Pi_R$; if moreover Axiom~I holds, then $\Pi_R|_R=\mathrm{id}_R$ so $r\in F_r$ for every $r\in R$, and $X=\bigsqcup_{r\in R}F_r$.
\end{definition}

The fiber decomposition is the set--theoretic quotient of $X$ under the equivalence relation $x\sim_\Pi y\iff\Pi_R(x)=\Pi_R(y)$; if Axiom~I holds, this quotient is canonically identified with $R$ via $r\leftrightarrow F_r$.

\begin{proposition}[Fiber--measure annihilation]\label{prop:fiber-annihilation}
Assume Axiom~I holds. Assume moreover that $\{r\}\in\mathcal A$ and $F_r\in\mathcal A$ for every $r\in R$, and that one of the following global regularity hypotheses holds:
\begin{itemize}[nosep]
\item[\emph{(R--fin)}] $R$ is finite; or
\item[\emph{(R--ctbl)}] $R$ is countable and $\mu$ is $\sigma$--additive on $\mathcal A$ (so $\mathcal A$ is assumed to be a $\sigma$--algebra).
\end{itemize}
If Axiom~III(b) (projection--measure invariance) holds, then
\[
\mu(F_r\setminus\{r\})=0\qquad(\forall\,r\in R),
\]
and consequently $\mu(X\setminus R)=0$: the measure $\mu$ is concentrated on the representative sector $R$.
\end{proposition}

\begin{proof}
By Axiom~I, $\Pi_R|_R=\mathrm{id}_R$, so $\Pi_R(r)=r$ and $r\in F_r=\Pi_R^{-1}(\{r\})$ for every $r\in R$. Hence
\[
F_r=\{r\}\sqcup(F_r\setminus\{r\})
\]
is a disjoint union of sets in $\mathcal A$. By invariance (Axiom~III(b)) applied to $B=\{r\}\subseteq R$,
\[
\mu(\Pi_R^{-1}(\{r\}))=\mu(\{r\}),
\]
so finite additivity gives
\[
\mu(\{r\})+\mu(F_r\setminus\{r\})=\mu(\{r\}),
\]
hence $\mu(F_r\setminus\{r\})=0$.

For the global claim $\mu(X\setminus R)=0$, decompose
\[
X\setminus R=\bigsqcup_{r\in R}(F_r\setminus\{r\}).
\]
Under (R--fin) this is a finite disjoint union of null sets, and finite additivity gives $\mu(X\setminus R)=0$. Under (R--ctbl) the union is countable, and $\sigma$--additivity of $\mu$ applied to the countable disjoint union gives again $\mu(X\setminus R)=0$.
\end{proof}

\begin{remark}[On the regularity hypothesis]\label{rem:regularity-hypothesis}
Without either hypothesis (R--fin) or (R--ctbl), the fiber--wise conclusion $\mu(F_r\setminus\{r\})=0$ still holds for each individual $r\in R$, but the global conclusion $\mu(X\setminus R)=0$ may fail, since finite additivity alone does not control infinitely many null sets. A finitely additive measure on an uncountable $R$ may assign positive total mass to an uncountable disjoint union of null sets. The hypotheses (R--fin), (R--ctbl) are the natural setting for the structural reduction theorem below; extensions to the uncountable $\sigma$--additive case require additional assumptions (e.g., separability of $\mathcal A$) which we do not pursue.
\end{remark}

\begin{remark}[Relation to Proposition~\ref{prop:muI-zero}]\label{rem:fiber-vs-muI}
Proposition~\ref{prop:fiber-annihilation} strengthens Proposition~\ref{prop:muI-zero} (support concentration under full partition): the full partition hypothesis $X=R\sqcup I$ is replaced by the weaker measurability condition on fibers. The conclusion $\mu(X\setminus R)=0$ now holds for any structural model whose fibers are measurable, regardless of whether $X=R\sqcup I$.
\end{remark}

\begin{theorem}[Quotient factorization]\label{thm:quotient-factorization}
Let
\[
\mathcal D=(X,\mathcal A,\mu,\mu^{\otimes2},R,I,\Pi_R,G,E_0,\eta)
\]
be a pre-structural datum (Definition~\ref{def:structural-datum}) with $\eta\in[0,1)$ and $\mu(X)<\infty$, satisfying the measurability hypothesis of Proposition~\ref{prop:fiber-annihilation} and one of the regularity hypotheses \emph{(R--fin)} or \emph{(R--ctbl)} of that proposition. Let $\mathcal D|_R$ denote the restricted datum on $R$ (with underlying set $R$, algebra $\mathcal A\cap\mathcal P(R)$, measure $\mu|_R$, restricted product charge, trivial retraction $\mathrm{id}_R$, restricted relation $G|_R:=G\cap(R\times R)$, scalar $\mu(R)$, and coupling parameter $\eta$).
\begin{enumerate}[label=(\roman*)]
\item (\emph{curvature load reduces to $R$}) If $\mathcal D$ satisfies Axiom~I and Axiom~III(b), then $\mu(X\setminus R)=0$, and for all $B\in\mathcal A$, writing $B_R:=B\cap R$,
\[
\mu^{\otimes2}\bigl((B\times X)\cap G\bigr)=\mu^{\otimes2}\bigl((B_R\times R)\cap (G|_R)\bigr).
\]
\item (\emph{coupling laws are equivalent}) If $\mathcal D$ satisfies Axiom~I and Axiom~III(b), then the coupling law (Axiom~III(c)) holds on $\mathcal D$ if and only if it holds on $\mathcal D|_R$.
\item (\emph{admissibility reduces}) If $\mathcal D$ satisfies Axioms~I, II, III(a), and~III(b) as well as the full partition condition $X=R\sqcup I$, then $\mathcal D$ is an admissible structural model (i.e., an object of $\mathbf{Struct}_\eta$) if and only if $\mathcal D|_R$ is an admissible structural model (i.e., an object of $\mathbf{Struct}^{\mathrm{id}}_\eta$). In particular, whenever the restricted core $\mathcal D|_R$ also satisfies the full hypotheses of Theorem~\ref{thm:pi-id-classification}, the classification of admissible structural models reduces to the identity--retraction classification of Theorem~\ref{thm:pi-id-classification}. (The hypothesis ``Axiom~III(b)'' is not redundant: the datum $\mathcal D_{\neg\mathrm{III}}$ of Proposition~\ref{prop:indep-III} satisfies Axioms~I, II, III(a), III(c), fiber measurability, (R--fin), and the full partition $X=R\sqcup I$, yet fails Axiom~III(b); its restriction $\mathcal D_{\neg\mathrm{III}}|_R$ is admissible while $\mathcal D_{\neg\mathrm{III}}$ itself is not, so the equivalence fails without~III(b).)
\end{enumerate}
\end{theorem}

\begin{proof}
(i) By Proposition~\ref{prop:fiber-annihilation}, which applies since $\mathcal D$ satisfies Axiom~I and Axiom~III(b) together with the measurability and regularity hypotheses, $\mu(X\setminus R)=0$. Since $\mu(X\setminus R)=0$, monotonicity and the rectangle rule give
\[
\mu^{\otimes2}(S)\;\le\;\mu^{\otimes2}\bigl((X\setminus R)\times X\bigr)=\mu(X\setminus R)\,\mu(X)=0
\]
for every measurable $S\subseteq(X\setminus R)\times X$. Similarly, $\mu^{\otimes2}(S')=0$ for every measurable $S'\subseteq X\times(X\setminus R)$.

Decomposing $B=B_R\sqcup(B\setminus R)$ with $B_R:=B\cap R$ (no full--partition hypothesis assumed) and using $R\sqcup(X\setminus R)=X$:
\[
(B\times X)\cap G
=\bigl((B_R\times R)\cap G\bigr)\sqcup\bigl((B_R\times(X\setminus R))\cap G\bigr)\sqcup\bigl(((B\setminus R)\times X)\cap G\bigr).
\]
The second term lies in $X\times(X\setminus R)$, hence has $\mu^{\otimes2}$--measure zero. The third term lies in $(X\setminus R)\times X$ (since $B\setminus R\subseteq X\setminus R$), hence also has $\mu^{\otimes2}$--measure zero. Therefore
\[
\mu^{\otimes2}((B\times X)\cap G)=\mu^{\otimes2}((B_R\times R)\cap G)=\mu^{\otimes2}((B_R\times R)\cap G|_R),
\]
where the last equality uses $(B_R\times R)\cap G=(B_R\times R)\cap(G\cap(R\times R))=(B_R\times R)\cap G|_R$, valid because $B_R\subseteq R$.

(ii) Assume first that the coupling law (Axiom~III(c)) holds on $\mathcal D$. For $B\in\mathcal A$ with $B_R=B\cap R$, by~(i) and Remark~\ref{rem:preimage} ($\Pi_R^{-1}(B)=\Pi_R^{-1}(B_R)$):
\[
\mu^{\otimes2}((B_R\times R)\cap G|_R)=\mu(B)+\eta\,\mu^{\otimes2}((\Pi_R^{-1}(B_R)\times R)\cap G|_R).
\]
Since $\mu(X\setminus R)=0$ we have $\mu(B)=\mu(B_R)$. Within $\mathcal D|_R$, the retraction is $\mathrm{id}_R$, so $\Pi_R^{-1}(B_R)\cap R=B_R$. Substituting:
\[
\mu^{\otimes2}((B_R\times R)\cap G|_R)=\mu(B_R)+\eta\,\mu^{\otimes2}((B_R\times R)\cap G|_R),
\]
which is the coupling law on $\mathcal D|_R$ applied to $B_R$. As $B$ ranges over $\mathcal A$, $B_R$ ranges over $\mathcal A\cap\mathcal P(R)$, so the coupling law on $\mathcal D|_R$ follows. The converse runs the same argument in reverse: given the coupling law on $\mathcal D|_R$, identity~(i) lifts it to $\mathcal D$.

(iii) We show the equivalence of admissibility. Axioms~I and~II for $\mathcal D$ imply Axioms~I and~II for $\mathcal D|_R$ (with retraction $\mathrm{id}_R$ and $G|_R$ inheriting reflexivity, symmetry, idempotence from $G$; see also the explicit argument for $G\circ G=G$ restricted to $R\times R$ given in the proof of Proposition~\ref{prop:restriction-functor}). For Axiom~III(a) on $\mathcal D|_R$: since $I\subseteq X\setminus R$ and $\mu(X\setminus R)=0$ by~(i), we have $\mu(I)=0$. Axiom~III(a) on $\mathcal D$ gives $\mu(R)+\mu(I)=E_0>0$, hence $\mu(R)=E_0>0$, which is precisely Axiom~III(a) on $\mathcal D|_R$ with scalar $\mu(R)$. Axiom~III(b) on $\mathcal D|_R$ is trivial because the retraction is $\mathrm{id}_R$. Axiom~III(c) on $\mathcal D|_R$ is equivalent to Axiom~III(c) on $\mathcal D$ by part~(ii). Combining, $\mathcal D$ satisfies Axioms~I--III if and only if $\mathcal D|_R$ does, i.e., $\mathcal D$ is an admissible structural model if and only if $\mathcal D|_R$ is. This proves the equivalence of admissibility. The final classification claim (reduction to the identity--retraction classification of Theorem~\ref{thm:pi-id-classification}) follows by applying Theorem~\ref{thm:pi-id-classification} to $\mathcal D|_R$ under its full hypotheses.
\end{proof}

\begin{corollary}[Closed form of the curvature load]\label{cor:curvature-load-closed-form}
Under the hypotheses of Theorem~\ref{thm:quotient-factorization}, assume moreover that $\mathcal D$ is admissible and that the restricted core $\mathcal D|_R$ satisfies the full hypotheses of Theorem~\ref{thm:pi-id-classification}. Then the observed curvature load on any measurable $B\subseteq X$ is given by
\[
\mu^{\otimes2}((B\times X)\cap G)=\sum_{k\in K}\mu(B_R\cap C_k)\,\mu(C_k),
\]
where $\{C_k\}_{k\in K}$ are the $G|_R$--equivalence classes of the identity--retraction core $\mathcal D|_R$, and $\mu(C_k)\in\{0,(1-\eta)^{-1}\}$ for each class, with $\mu(C_k)=(1-\eta)^{-1}$ on positive--mass classes.
\end{corollary}

\begin{proof}
Combine Theorem~\ref{thm:quotient-factorization}(i) with Theorem~\ref{thm:pi-id-classification}, applied to the admissible core $\mathcal D|_R$ under its full hypotheses.
\end{proof}

\begin{remark}[Structure of the reduction]\label{rem:reduction-structure}
\sloppy
Theorem~\ref{thm:quotient-factorization}(iii) exhibits the pair $(\mathbf{Struct}^{\mathrm{id}}_{\eta},\allowbreak\mathbf{Struct}_{\eta})$ as analogous to a space and a null--set extension of it: passing from $\mathcal D|_R$ to $\mathcal D$ adds only measure--zero fiber data, invisible to the coupling law and the curvature load. This parallels the measure--theoretic fact that two measures agreeing off a null set are operationally indistinguishable.
\end{remark}

\begin{remark}[Place of the quotient factorization in the paper]\label{rem:place-quotient}
\sloppy
Theorem~\ref{thm:quotient-factorization} bridges three earlier structural results into a single statement:
\begin{itemize}[nosep]
\item It generalizes Proposition~\ref{prop:muI-zero} (support concentration), extending $\mu(I)=0$ from the full partition case to the fiber--measurability case.
\item It promotes the restriction functor $\rho$ of Proposition~\ref{prop:restriction-functor} from a categorical operation to a structural equivalence on curvature loads.
\item It reduces the general classification problem to Theorem~\ref{thm:pi-id-classification}, under measurability hypotheses on fibers and the full hypotheses of Theorem~\ref{thm:pi-id-classification} on the restricted core.
\end{itemize}
\end{remark}

% -------------------------------------------------------

\section{Examples}
We conclude with two finite examples.
\subsection{Finite representative/hidden pattern}
\begin{example}[Finite representative/hidden pattern]
\sloppy
Let $X$ be a finite nonempty set, partitioned into nonempty blocks
$(C_j)_{j\in J}$.
For each $j$, choose a distinguished element $r_j\in C_j$ and define
\[
R := \{r_j : j\in J\},\qquad I := X\setminus R.
\]
One may view $R$ as the observable representative sector and $I$ as a hidden sector; we define $\Pi_R:X\to R$ by $\Pi_R(x)=r_j$ whenever $x\in C_j$.
Let $\mathcal A=\mathcal P(X)$ and choose weights $w_j\in\{0,1\}$ with $\sum_j w_j>0$.
Define $\mu(B)=\sum_{j : r_j\in B} w_j$, and let
\[
G := \bigcup_{j\in J} (C_j\times C_j)
\]
be the union of equivalence relations on each block.
Set $E_0=\mu(X)$ and $\eta=0$.
Then Proposition~\ref{prop:finite-class-model} shows that
\[
\mathcal M=(X,\mathcal A,\mu,\mu^{\otimes2},R,I,\Pi_R,G,E_0,\eta)
\]
is a structural model.
Each block $C_j$ represents a cluster of points sharing the same hidden content, with $r_j$ as the open representative.
\end{example}
\subsection{Abstract equivalence/bridge pattern}
\begin{example}[Abstract equivalence/bridge pattern]
\sloppy
Let $X$ be any nonempty set equipped with an equivalence relation $\approx$.
Assume $X/\!\approx$ is finite, and choose a representative $r_j$ for each class $C_j$.
Set
\[
R := \{r_j\},\qquad I := X\setminus R,\qquad
G := \bigcup_j (C_j\times C_j),
\]
and let $\Pi_R:X\to R$ send each $x\in C_j$ to $r_j$.
Let $\mathcal A=\mathcal P(X)$.
Choose weights $w_j\in\{0,1\}$ with $\sum_j w_j>0$ and define
\[
\mu(B) = \sum_{j : r_j\in B} w_j.
\]
Then by Proposition~\ref{prop:finite-class-model}, the tuple
$\mathcal M=(X,\mathcal A,\mu,\mu^{\otimes2},R,I,\Pi_R,G,E_0,\eta)$ with $E_0=\mu(X)$ and $\eta=0$
is a structural model.
Heuristically, $G$ encodes a bridge relation between configurations, grouping those
that share the same structural type.
The set $R$ collects canonical representatives of these types, and $\Pi_R$ retracts the
full space $X$ onto the simpler representative space $R$ while preserving the chosen
measure on $R$.
The axioms formalize this pattern in a ZFC--internal way.
\end{example}
% -------------------------------------------------------
\section{Conclusion}

We have set up a minimal ZFC--internal axiom system for admissible structural models — pre-structural data $(X,\mathcal A,\mu,\mu^{\otimes2},R,I,\Pi_R,G,E_0,\eta)$ satisfying Axioms~I, II, and~III — and studied its basic properties.

The axiom system is consistent in ZFC (Section~\ref{sec:models}), and its three axioms, together with the three subclauses of Axiom~III, are mutually independent (Theorem~\ref{thm:independence-main}, Corollary~\ref{cor:minimality}). The coupling law of Axiom~III admits an operator--theoretic reformulation as the unique bounded finitely additive solution of the Banach--contraction equation $f=T_\eta f$, with closed form $f_*(B)=\mu(B)+\tfrac{\eta}{1-\eta}\mu(\Pi_R^{-1}(B))$ (Theorem~\ref{thm:coupling-fixed-point}). Admissible structural models with a common $\eta$ form a category $\mathbf{Struct}_\eta$ (Section~\ref{sec:category}); under fiber measurability and either (R--fin) or (R--ctbl), the general admissibility problem reduces to the identity--retraction case (Theorem~\ref{thm:quotient-factorization}), where the classification of Theorem~\ref{thm:pi-id-classification} yields the block dichotomy $\mu(C_k)\in\{0,(1-\eta)^{-1}\}$. Under $\sigma$--additive probability normalization, the $B=X$ instance of the coupling law forces $\eta=0$ (Section~\ref{sec:sigma-additive}), so $\eta\neq 0$ families live in the finitely additive setting or among non--probability $\sigma$--additive measures.

\paragraph{Further directions.}
Several extensions suggest themselves: replacing $\mathcal A$ by a $\sigma$--algebra under hypotheses giving a well--defined product charge; enriching $G$ to a weighted kernel $K:X\times X\to[0,\infty)$ and comparing with graph limits \cite{Lovasz}; further categorical development, in particular an adjoint characterization of the restriction $\rho:\mathbf{Struct}_\eta^{\mathrm{fib}}\to\mathbf{Struct}^{\mathrm{id}}_\eta$ \cite{MacLane,Riehl}; operator--theoretic variants replacing $\mu$ by a state or trace; and scaling limits $X_n\to X$ controlled by local instances of the coupling law. Removing the fiber hypotheses in the quotient--factorization reduction, for instance in the uncountable $\sigma$--additive case, also requires further assumptions such as separability of $\mathcal A$.

% -------------------------------------------------------

\begin{thebibliography}{99}

\bibitem{Jech}
T.~Jech, \emph{Set Theory}, 3rd ed., Springer, 2003.

\bibitem{Kunen}
K.~Kunen, \emph{Set Theory}, North--Holland, 1980.

\bibitem{Enderton}
H.~Enderton, \emph{Elements of Set Theory}, Academic Press, 1977.

\bibitem{Levy}
A.~Levy, \emph{Basic Set Theory}, Springer, 1979.

\bibitem{Halmos}
P.~Halmos, \emph{Measure Theory}, Springer, 1974.

\bibitem{Bogachev}
V.~Bogachev, \emph{Measure Theory}, 2 vols., Springer, 2007.

\bibitem{Rao}
K.~P.~S.~Bhaskara Rao and M.~Bhaskara Rao, \emph{Theory of Charges}, Academic Press, 1983.

\bibitem{Fremlin}
D.~H.~Fremlin, \emph{Measure Theory}, 5 vols., Torres Fremlin, Colchester, 2000--2008.

\bibitem{AliprantisBorder}
C.~D.~Aliprantis and K.~C.~Border, \emph{Infinite Dimensional Analysis: A Hitchhiker's Guide}, 3rd ed., Springer, 2006.

\bibitem{Kuratowski}
K.~Kuratowski, \emph{Topology I}, Academic Press, 1966.

\bibitem{Ore}
O.~Ore, ``Theory of Equivalence Relations,'' \emph{Duke Math. J.} 9 (1942), 573--627.

\bibitem{Sabidussi}
G.~Sabidussi, ``Graph Derivatives,'' \emph{Math. Z.} 76 (1961), 385--401.

\bibitem{KolokoltsovMaslov}
V.~N.~Kolokoltsov and V.~P.~Maslov, \emph{Idempotent Analysis and Its Applications}, Math. Appl. 401, Kluwer Academic, 1997.

\bibitem{BauschkeCombettes}
H.~H.~Bauschke and P.~L.~Combettes, \emph{Convex Analysis and Monotone Operator Theory in Hilbert Spaces}, 2nd ed., Springer, 2017.

\bibitem{Lovasz}
L.~Lov\'asz, \emph{Large Networks and Graph Limits}, American Mathematical Society Colloquium Publications 60, AMS, 2012.

\bibitem{KrantzEtAl}
D.~H.~Krantz, R.~D.~Luce, P.~Suppes, and A.~Tversky, \emph{Foundations of Measurement}, 3 vols., Academic Press, 1971--1990.

\bibitem{MacLane}
S.~Mac~Lane, \emph{Categories for the Working Mathematician}, 2nd ed., Springer, 1998.

\bibitem{Awodey}
S.~Awodey, \emph{Category Theory}, 2nd ed., Oxford Univ. Press, 2010.

\bibitem{Borceux}
F.~Borceux, \emph{Handbook of Categorical Algebra}, 3 vols., Cambridge Univ. Press, 1994.

\bibitem{Riehl}
E.~Riehl, \emph{Category Theory in Context}, Dover, 2016.

\end{thebibliography}

% -------------------------------------------------------

\end{document}